\chapter{Arithmetic Chern--Simons and boundary realization}
\label{v4-arith:cs:chapter}

A gauge theory assigns an amplitude to a field and sums over fields when
regions are glued. These are two different operations. For the group
\(C_2=\{1,-1\}\), the sign representation has Wilson values \(1\) and \(-1\)
on the two circle holonomies. With trivial action their normalized average
is zero. Specifying one holonomy instead gives one of the two values. An
arithmetic Wilson comparison must determine which operation produces the
Frobenius trace.

Finite gauge groups make the distinction exact. Automorphisms determine the
integration measure, a three-cocycle determines the action, and
transgression determines the boundary line. Over a number field, arithmetic
duality supplies a different integration map in cohomological degree three.
It gives a scalar Chern--Simons invariant and torsors over local
representation spaces. In the abelian case it also gives a linking pairing
whose Gaussian sum computes Wilson correlations. These are constructions on
specified arithmetic objects.

We develop these finite theories before passing to congruence towers. The
transition maps expose the lifting measures required by a profinite limit. A
free holomorphic--topological model then supplies an explicit physical
boundary comparison, with its current algebra, coefficient complex and
Gaussian cap. The final question is to construct arithmetic bulk and
boundary objects for which these operations agree. Neither an Euler product
nor a boundary character supplies those maps.

The prerequisites are group cohomology, finite groupoids and the coefficient
and trace distinctions of Chapters~\ref{v4-arith:coefficients:chapter}
and~\ref{v4-arith:traces:chapter}. Continuous cochains are used for
profinite groups. Tensor products are algebraic unless a completion is
specified.

\section{Finite gauge fields and their integration}
\label{v4-arith:cs:finite}

A \emph{finite groupoid} is a groupoid with finitely many isomorphism
classes and finite automorphism groups. Its integral of an
isomorphism-invariant complex function is
\begin{equation}
 \int_{\mathcal X} f
 =\sum_{[x]\in\pi_0\mathcal X}
       \frac{f(x)}{|\operatorname{Aut}_{\mathcal X}(x)|},
 \qquad
 |\mathcal X|=\int_{\mathcal X}1.
 \label{v4-arith:cs:eq:integral}
\end{equation}
The symbol \(BG\) denotes the one-object groupoid with automorphism group
\(G\). Thus \(|BG|=1/|G|\). For a finite \(G\)-set \(X\), the action
groupoid \(X//G\) has objects \(X\) and arrows the action of \(G\).
Orbit--stabilizer gives
\begin{equation}
 \int_{X//G}f=\frac1{|G|}\sum_{x\in X}f(x).
 \label{v4-arith:cs:eq:action-integral}
\end{equation}

For a connected space with fundamental group \(\Gamma\), finite gauge fields
with group \(G\) form \(\mathcal F_G=\operatorname{Hom}(\Gamma,G)//G\),
whenever this groupoid is finite. An object is a homomorphism \(\sigma\).
Its automorphism group is the centralizer of \(\sigma(\Gamma)\). In the
arithmetic case \(\Gamma\) is a specified profinite fundamental group and
the homomorphisms are continuous. Ramification conditions are part of that
specification. They cannot be recovered from the finite group alone.

\subsection{Integration over homotopy fibers}

Let \(p:\mathcal X\to\mathcal Y\) be a functor of finite groupoids. Its
\emph{homotopy fiber} at \(y\) has objects \((x,\phi)\), where
\(\phi:p(x)\xrightarrow{\sim}y\). An arrow \(a:(x,\phi)\to(x',\phi')\)
satisfies \(\phi' p(a)=\phi\). Define
\[
 (p_!f)(y)=\int_{\mathcal X_y}f(x),\qquad
 (p^*g)(x)=g(p(x)).
\]
Keeping \(\phi\) is necessary for the measure.

\begin{theorem}[Finite integration]
\label{v4-arith:cs:thm:fubini}
For functors of finite groupoids the following identities hold:
\[
 \int_{\mathcal Y}p_!f=\int_{\mathcal X}f,
 \qquad (q p)_!=q_!p_!,
 \qquad p_!(f\,p^*g)=(p_!f)g.
\]
In a homotopy Cartesian square
\[
\begin{tikzcd}
 \mathcal X' \ar[r,"q'"]\ar[d,"p'"']&
 \mathcal X\ar[d,"p"]\\
 \mathcal Y'\ar[r,"q"']&\mathcal Y,
\end{tikzcd}
\] one has \(q^*p_!=p'_!q'^*\).
\end{theorem}

\begin{proof}
Fix an isomorphism class \([x]\) and a representative \(y\) of \([p(x)]\).
The set \(\operatorname{Iso}(p(x),y)\) has \(|\operatorname{Aut}y|\)
elements. The contribution of the component over \([x]\) to \(\mathcal X_y\)
is its action groupoid under \(\operatorname{Aut}x\), and hence has mass
\(|\operatorname{Aut}y|/|\operatorname{Aut}x|\). Integration over \([y]\)
cancels the numerator and gives \eqref{v4-arith:cs:eq:integral}.

For a composite, an object in an iterated fiber consists of \(x,y\), an
isomorphism \(p(x)\to y\), and an isomorphism \(q(y)\to z\). Composing the
latter with the image of the former identifies the iterated fiber with the
fiber of \(qp\), with its automorphisms. Applying the same count proves
composition. The projection formula follows because \(g(p(x))=g(y)\) in the
fiber. Finally, the fiber of \(p'\) at \(y'\) is equivalent to \(\mathcal
X_{q(y')}\). This equivalence preserves the integrand and proves base
change.
\end{proof}

For example, \(\mathrm{pt}\times^h_{BG}\mathrm{pt}\) is the discrete set
\(G\). The two trivial objects must be supplied with an isomorphism, and
that is an element of \(G\). Its mass is \(|G|\). The pullback of the two
sets of isomorphism classes is one point and has mass one. It gives the
wrong gluing factor as soon as \(G\) is nontrivial.

\subsection{Boundary lines and adjoints}

A Hermitian line \(L\) on a groupoid \(\mathcal Y\) assigns a Hermitian
one-dimensional space to each object and unitary transport to each arrow.
The transport is functorial. Its invariant sections form a Hilbert space
\begin{equation}
 H(\mathcal Y,L)=\Gamma(\mathcal Y,L),\qquad
 \langle s,t\rangle
 =\sum_{[y]}\frac{\langle s(y),t(y)\rangle_{L_y}}
                         {|\operatorname{Aut}y|}.
 \label{v4-arith:cs:eq:boundary-hilbert}
\end{equation}
The pairing is linear in its first variable. If an automorphism acts
nontrivially on \(L_y\), every invariant section vanishes at \(y\). The
corresponding component contributes no state.

Pullback of sections along \(p\) has an adjoint. For \(s\in H(\mathcal
X,p^*L)\), it is
\begin{equation}
 (p_!s)(y)=
 \sum_{[(x,\phi)]\in\pi_0\mathcal X_y}
  \frac{L(\phi)s(x)}{|\operatorname{Aut}(x,\phi)|}.
 \label{v4-arith:cs:eq:line-push}
\end{equation}
The transport makes each summand lie in \(L_y\). Fubini and unitarity give
\(\langle p_!s,t\rangle=\langle s,p^*t\rangle\). Thus orientation reversal,
when it conjugates the action and reverses the field span, reverses the
corresponding operator by its Hilbert adjoint.

This integral is distinct from categorical pushforward of a representation.
For \(p:BG\to\mathrm{pt}\), the invariants of the trivial complex
representation are \(\mathbb C\), of dimension one. The integral of the
constant function is \(1/|G|\). Both operations exist. Identifying one with
a trace of the other requires a specified trace construction and
normalization.

\subsection{Circle and surface counts}

For the circle, \(\mathcal F_G(S^1)=G//G\), with conjugation action, so its
mass is one. The space of class functions on this groupoid has dimension the
number of conjugacy classes. These two numbers already distinguish a
partition mass from a state-space dimension.

For a closed oriented genus-\(g\) surface, the untwisted mass is
\begin{equation}
 |\mathcal F_G(\Sigma_g)|
 =\frac1{|G|}\#\left\{(a_i,b_i)_{i=1}^g:
              \prod_i[a_i,b_i]=1\right\}
 =|G|^{2g-2}\sum_{\rho\in\widehat G}
                      (\dim\rho)^{2-2g}.
 \label{v4-arith:cs:eq:genus}
\end{equation}
Here \(\widehat G\) is the set of irreducible complex representations.

To verify the formula, put \(K=\sum_{a,b\in G}[a,b]\in\mathbb C[G]\). For an
irreducible representation of dimension \(d\), averaging
\(\rho(a)\rho(b)\rho(a)^{-1}\) over \(a\) gives \(|G|\chi_\rho(b)I/d\).
Character orthogonality then gives \(\rho(K)=|G|^2I/d^2\). The coefficient
of the identity in \(K^g\) is \(|G|^{-1}\) times its regular trace. Since
\(\rho\) occurs \(d\) times in the regular representation, division by the
additional \(|G|\) gives \eqref{v4-arith:cs:eq:genus}.

For \(S_3\), whose irreducible dimensions are \(1,1,2\), the torus and
genus-two masses are respectively \(3\) and \(81\). In three dimensions a
different count occurs:
\[
 |\mathcal F_G(T^3)|
 =\frac{\#\{(a,b,c):a,b,c\text{ commute pairwise}\}}{|G|}
 =\sum_{[a]}\#\widehat{C_G(a)}.
\]
For fixed \(a\), the number of commuting pairs in its centralizer is
\(|C_G(a)|\) times the number of its conjugacy classes. Summing over the
class of \(a\) proves the equality. For \(S_3\) the three centralizers are
\(S_3,C_2,C_3\), and the answer is \(3+2+3=8\).

\section{Three-cocycles and finite topological gluing}
\label{v4-arith:cs:topological}

Let \(G\) be finite. A normalized multiplicative three-cocycle
\(\omega:G^3\to U(1)\) is one when any argument is the identity and
satisfies
\begin{equation}
 \omega(h,k,l)\omega(g,hk,l)\omega(g,h,k)
 =\omega(gh,k,l)\omega(g,h,kl).
 \label{v4-arith:cs:eq:cocycle}
\end{equation}
For an oriented triangulated closed three-manifold, assign \(g_{ij}\in G\)
to oriented edges, with \(g_{ji}=g_{ij}^{-1}\) and \(g_{ij}g_{jk}=g_{ik}\)
on each face. The action phase is
\[
 a(g)=\prod_{[0123]}
           \omega(g_{01},g_{12},g_{23})^{\epsilon_{0123}},
 \qquad
 Z_\omega(M)=|G|^{-N_0}\sum_{\text{flat }g}a(g),
\]
where \(N_0\) is the number of vertices and \(\epsilon_{0123}\) records
orientation. The factor \(|G|^{-N_0}\) integrates the action of vertex gauge
transformations. It remains necessary when the action has stabilizers.

The two-to-three Pachner identity is precisely
\eqref{v4-arith:cs:eq:cocycle}. In a one-to-four move, an extension of a
fixed flat boundary field has \(|G|\) choices for the new vertex label. The
cocycle identity makes each extended phase equal to the original phase, and
the new vertex factor cancels these choices. Normalization handles
degenerate simplices. This gives the finite Dijkgraaf--Witten state sum.
Boundary amplitudes require the line determined by the same cocycle.

\subsection{The degree of transgression}

Suppose a space \(K\) and its field space admit an evaluation map
\(\mathrm{ev}:K\times\mathcal F_G(K)\to BG\). For a chain \(c\in C_d(K)\),
define the transgression of an additive cochain \(u\) by
\((\tau_cu)(a)=\mathrm{ev}^*u(c\times a)\). The cross-product boundary
formula gives
\begin{equation}
 \delta\tau_cu=(-1)^d
                  \bigl(\tau_c\delta u-\tau_{\partial c}u\bigr).
 \label{v4-arith:cs:eq:transgression}
\end{equation}
Indeed, evaluating \(\delta u\) on \(c\times a\) gives \(u(\partial c\times
a)+(-1)^du(c\times\partial a)\). Solving for the second term proves the
identity.

For a degree-three cocycle, integration over a closed three-manifold gives a
degree-zero phase. Over a closed surface it gives a degree-one cocycle,
hence a flat line on the field groupoid. Over a circle it gives a degree-two
cocycle, hence twisting data for representations of that field groupoid. At
a point the degree-three class supplies an associator.

The degrees are
\[
\begin{array}{c|c|c}
 \text{dimension integrated}&\text{degree remaining}&\text{recipient}\\\hline
 3&0&\text{action phase}\\
 2&1&\text{line on fields}\\
 1&2&\text{twisted representation category}\\
 0&3&\text{monoidal associator}.
\end{array}
\]

In particular, the boundary of a tubular neighborhood of a line in a
three-manifold is a surface. Its transgression is a line on surface fields.
The twisted category belongs one categorical degree lower, at a circle, with
the indicated transgression.

For a three-chain with boundary, \eqref{v4-arith:cs:eq:transgression} says
that the boundary of its action cochain is the boundary-line cocycle. Its
exponential is therefore a section of that line, rather than an invariant
scalar on unrestricted boundary fields. Changing the orientation conjugates
the line and the phase. These constructions are the finite topological ones
of Freed--Quinn, Theorems~1.7 and~2.13. They use oriented manifolds and the
actual integration of their chains. An arithmetic analogue must supply its
own duality and boundary integration.

\begin{theorem}[Finite gluing on specified field data]
\label{v4-arith:cs:thm:gluing}
Suppose regions have finite field groupoids, boundary Hermitian lines, and
invariant action sections. Assume gluing regions gives the homotopy pullback
of their field groupoids, and that the action sections contract under the
boundary-line pairing.

Assign to a boundary its Hilbert space
\eqref{v4-arith:cs:eq:boundary-hilbert}. For a region from \(\Sigma_0\) to
\(\Sigma_1\), let \(r_i\) be restriction of fields and let \(a\) be the
section of \(r_0^*L_0^*\otimes r_1^*L_1\). Then
\begin{equation}
 T_Ms=(r_1)_!\bigl(a\,r_0^*s\bigr)
 \label{v4-arith:cs:eq:region-operator}
\end{equation}
composes under gluing.

A cylinder whose field restriction is the identity and whose section is the
identity gives the identity operator. Disjoint unions give tensor products.
Reversing a region and conjugating its action gives \(T_M^*\).
\end{theorem}

\begin{proof}
In the composite of two expressions \eqref{v4-arith:cs:eq:region-operator},
base change replaces restriction followed by integration with integration
over the homotopy pullback. The two action sections multiply and contract
there to the glued section. Fubini combines the two integrations. The
cylinder and disjoint-union assertions follow from the identity functor and
multiplication of automorphism orders. Finally,
\eqref{v4-arith:cs:eq:line-push} is adjoint to pullback, and the adjoint of
multiplication by a line section uses its Hermitian conjugate. This reverses
the span and proves the last assertion.
\end{proof}

Descent of fields supplies the homotopy pullback in this theorem. The action
sections and their contraction are additional hypotheses. An equality of
field groupoids without those sections does not determine the glued
amplitude.

\section{The arithmetic action and its boundary torsor}
\label{v4-arith:cs:arithmetic-action}

Let \(F\) be a totally imaginary number field containing \(\mu_n\), put
\(X=\operatorname{Spec}\mathcal O_F\), and fix \(\zeta:\mathbb
Z/n\xrightarrow{\sim}\mu_n\). Arithmetic duality gives an invariant map
\[
 \operatorname{inv}_\zeta:H^3_{\mathrm{\acute et}}(X,\mathbb Z/n)
       \xrightarrow{\sim}\tfrac1n\mathbb Z/\mathbb Z.
\]
The choice of \(\zeta\) identifies the constant coefficients with the
roots-of-unity coefficients in this map.

Write \(\pi=\pi_1^{\mathrm{\acute et}}(X)\). For a finite group \(G\) and
\(c\in Z^3(G,\mathbb Z/n)\), define
\begin{equation}
 \operatorname{CS}_c(\sigma)
 =\operatorname{inv}_\zeta\bigl(j^3\sigma^*[c]\bigr),
 \qquad
 j^3:H^3_{\mathrm{cont}}(\pi,\mathbb Z/n)
                \longrightarrow H^3_{\mathrm{\acute et}}(X,\mathbb Z/n).
 \label{v4-arith:cs:eq:kim}
\end{equation}
This is Kim's finite arithmetic action. The displayed comparison \(j^3\) is
used as a map. No equivalence between the entire group cohomology and
\'{e}tale cohomology is needed.

The value in \eqref{v4-arith:cs:eq:kim} is a scalar in \(\tfrac1n\mathbb
Z/\mathbb Z\). Before evaluation, the input is a degree-three class. The
coefficient sheaf \(\mathbb G_m\) appearing in arithmetic duality is not the
finite gauge group \(G\). Nor does cohomological degree three make the
one-dimensional scheme \(X\) a real three-manifold.

Conjugation acts trivially on the pulled-back cohomology class, so the
action descends to the finite field groupoid
\(\operatorname{Hom}_{\mathrm{cont}}(\pi,G)//G\). Its finite partition
value, with the measure already fixed, is
\begin{equation}
 Z_{G,c}(X)=\sum_{[\sigma]}
 \frac{\exp(2\pi i\operatorname{CS}_c(\sigma))}
      {|C_G(\sigma(\pi))|}.
 \label{v4-arith:cs:eq:arithmetic-sum}
\end{equation}
Finiteness here is the arithmetic finiteness of the specified representation
set, as follows from bounded-degree extension finiteness with the stated
ramification conditions. A different arithmetic region needs its own such
assertion.

\subsection{An explicit gauge homotopy}

The boundary construction needs conjugation on cochains, including its
composition. Use the inhomogeneous differential with trivial \(G\)-action on
\(\mathbb Z/n\). For a two-cochain \(b\),
\[
 (db)(g,h,k)=b(h,k)-b(gh,k)+b(g,hk)-b(g,h).
\]

Put \(\operatorname{Ad}_a g=aga^{-1}\). For a normalized three-cocycle
\(c\), define
\begin{align}
 h_a(g,h)={}&c(a^{-1},\operatorname{Ad}_ag,\operatorname{Ad}_ah)
       -c(g,a^{-1},\operatorname{Ad}_ah)
       +c(g,h,a^{-1}),
 \label{v4-arith:cs:eq:prism}\\
 H_{a,b}(g)={}&c(b^{-1},a^{-1},\operatorname{Ad}_{ab}g)
       -c(b^{-1},\operatorname{Ad}_bg,a^{-1})
       +c(g,b^{-1},a^{-1}).\notag
\end{align}
The cochain identities are
\begin{equation}
 dh_a=c\circ\operatorname{Ad}_a-c,
 \qquad
 h_{ab}-h_a\circ\operatorname{Ad}_b-h_b=dH_{a,b}.
 \label{v4-arith:cs:eq:prism-identities}
\end{equation}

They follow by expanding the differential and cancelling the three-cocycle
identities on the successive simplices of a prism. The alternating signs in
\eqref{v4-arith:cs:eq:prism} are essential. These are also the explicit
conjugation homotopies in Chung--Kim--Kim--Park--Yoo, Appendix~B, in the
inhomogeneous convention above.

If \(db=\sigma^*c\), conjugation by \(a\) sends the primitive to
\(b+\sigma^*h_a\), a primitive for \((\operatorname{Ad}_a\sigma)^*c\). The
second identity in \eqref{v4-arith:cs:eq:prism-identities} makes these maps
compose after quotienting primitives by two-coboundaries. Thus the
construction retains the gauge action on the torsor itself.

\subsection{Local torsors and the global section}

Let \(S\) be a nonempty finite set of finite places containing every place
dividing \(n\). Put
\[
 U=\operatorname{Spec}\mathcal O_F[1/S],\qquad
 \pi_S=\pi_1^{\mathrm{\acute et}}(U).
\]

Choose
restriction maps from \(\pi_v=\operatorname{Gal}(\overline F_v/F_v)\) to
\(\pi_S\). For a local field
\(\sigma_v:\pi_v\to G\), consider
\begin{equation}
 P_v(\sigma_v)=
 \{b_v\in C^2(\pi_v,\mathbb Z/n):db_v=\sigma_v^*c\}/B^2.
 \label{v4-arith:cs:eq:local-torsor}
\end{equation}
Local cohomological dimension gives \(H^3(\pi_v,\mathbb Z/n)=0\), so this
set is nonempty. The difference of two primitives is a two-cocycle.
Consequently \(P_v(\sigma_v)\) is a torsor for \(H^2(\pi_v,\mathbb Z/n)\),
identified by local duality and \(\zeta\) with \(\tfrac1n\mathbb Z/\mathbb
Z\).

Take the product of these torsors and push it out along the sum of the local
invariant maps. The result is a \(\tfrac1n\mathbb Z/\mathbb Z\)-torsor
\(P_S(\sigma_S)\). Here the pushout of a torsor means its contracted product
along the specified homomorphism of acting groups. Exponentiation along
\(t\mapsto e^{2\pi it}\) makes the associated Hermitian complex line.
Formula~\eqref{v4-arith:cs:eq:prism} defines its maps under local
conjugations.

For a global \(\sigma:\pi_S\to G\), the same cohomological dimension theorem
gives a primitive \(b_+\) with \(db_+=\sigma^*c\). Its local restrictions
determine an element
\begin{equation}
 \operatorname{CS}_{U,c}(\sigma)
       =[\operatorname{res}_v b_+]_{v\in S}
       \quad\text{in }P_S(\operatorname{res}\sigma).
 \label{v4-arith:cs:eq:boundary-section}
\end{equation}
Changing \(b_+\) by a global two-cocycle does not change this element:
global reciprocity says that the sum of its local invariants is zero.
Changing it by a coboundary already has no effect in
\eqref{v4-arith:cs:eq:local-torsor}. The gauge homotopies make
\eqref{v4-arith:cs:eq:boundary-section} an invariant section over the
restriction functor from global to local fields.

The arithmetic input is vanishing in degree three, local invariants in
degree two, and reciprocity for restrictions of global degree-two classes.
It supplies the torsor and its section. Gluing arbitrary arithmetic bordisms
would add the corresponding regions, field descent and compatible
action-section identifications of Theorem~\ref{v4-arith:cs:thm:gluing}.

\subsection{The decomposition formula}

Suppose now that \(\sigma\) is unramified on \(X\). At each \(v\in S\), its
local restriction factors through \(\pi_v/I_v\cong\widehat{\mathbb Z}\).
This group has cohomological dimension one for the finite coefficients in
use. Thus its degree-three cocycle has a primitive \(b_{-,v}\) unique modulo
two-coboundaries. Inflation gives a canonical class of unramified local
primitives in \eqref{v4-arith:cs:eq:local-torsor}.

\begin{theorem}[Arithmetic decomposition]
\label{v4-arith:cs:thm:decomposition}
With the coefficients and hypotheses above,
\begin{equation}
 \operatorname{CS}_c(\sigma)
  =\sum_{v\in S}\operatorname{inv}_v
          [b_{-,v}-\operatorname{res}_v b_+].
 \label{v4-arith:cs:eq:decomposition}
\end{equation}
Each bracket is a two-cocycle class. The right side is independent of the
chosen global primitive and of the representatives of the local primitive
classes.
\end{theorem}

\begin{proof}
Use the compact-support mapping fiber, with differential
\[
 C_c^j=C^j(\pi_S,\mathbb Z/n)
             \oplus\prod_{v\in S}C^{j-1}(\pi_v,\mathbb Z/n),
 \qquad
 d(a,b)=(da,\operatorname{res}a-db).
\]
The pair \((\sigma^*c,b_-)\) represents the unramified class on \(X\).
Subtracting \(d(b_+,0)\) gives \((0,b_--\operatorname{res}b_+)\). Arithmetic
duality identifies its invariant with the sum of the displayed local
invariants. Reciprocity removes a change in \(b_+\), and local coboundaries
remove a change in \(b_-\). This is the mapping-fiber proof of
Chung--Kim--Kim--Park--Yoo, Theorem~1.1 and \S4.
\end{proof}

Both primitive classes belong to the same boundary torsor. Their difference
is a scalar. Reversing the difference reverses the action, so the
compact-support differential fixes a necessary sign.

\subsection{Compact \(p\)-adic coefficients}

There is also a torsor construction with compact coefficients. For an odd
prime \(p\), write \(\Gamma=\mathbb Z_p^\times=\Delta\times\Gamma_1\), where
\(\Delta\) has order \(p-1\) and \(\Gamma_1=1+p\mathbb Z_p\). Let \(A\) be a
\(p\)-adic Lie group with a homomorphism \(t:A\to\Gamma\), and let \(\chi\)
be the cyclotomic character. Fix an actual continuous automorphism
\(\sigma\) of \(\Gamma\) and restrict to representations \(\rho\) satisfying
\(t\rho=\sigma\chi\). Gauge transformations are taken in \(\ker t\).

Let \(\mathbb Z_p[[\Gamma]]_\rho\) have the left action through \(t\rho\). A
compatible choice of roots of unity gives a basis \(\zeta\) of \(\mathbb
Z_p(1)\). There is an equivariant isomorphism
\begin{equation}
 \mathbb Z_p[[\Gamma]]_\rho
   \longrightarrow
 \Lambda=\mathbb Z_p[[\Gamma]]\widehat\otimes\mathbb Z_p(1),
 \qquad
 [\gamma]\longmapsto
 [\sigma^{-1}\gamma]\otimes(\sigma^{-1}\gamma)\zeta.
 \label{v4-arith:cs:eq:padic-coefficients}
\end{equation}
The first factor on the right has left action through \(\chi\), and the
second has the Tate action. To check equivariance, put \(u=\chi(g)\). The
image of \([\sigma(u)\gamma]\) is \([u\sigma^{-1}\gamma]\otimes
u(\sigma^{-1}\gamma)\zeta\), which is exactly the action of \(g\) on the
displayed image.

For a continuous class \(c\in H^3(A,\mathbb Z_p[[\Gamma]])\), use
\eqref{v4-arith:cs:eq:padic-coefficients} on the pulled-back cocycle. Under
vanishing of degree-three continuous cohomology with these compact
coefficients, its local primitives form \(H^2(\pi_v,\Lambda)\)-torsors. The
augmentation \(\Lambda\to\mathbb Z_p(1)\), followed by the local invariant
\(H^2(\pi_v,\mathbb Z_p(1))\cong\mathbb Z_p\), pushes their product to a
\(\mathbb Z_p\)-torsor. A global primitive gives a section by the same
reciprocity argument. Gauge maps use the conjugation prism with this
coefficient action.

A sufficient setting uses finite quotient modules that compute the
continuous compact-coefficient cochains by an inverse sequence. Assume
finite degree-two cohomology and vanishing degree-three cohomology. The
finite degree-two systems satisfy the Mittag--Leffler condition. The derived
inverse-limit exact sequence then gives degree-three vanishing. The
augmentations, local duality and reciprocity must be compatible with this
sequence. These hypotheses describe the coefficient passage used in the
\(p\)-adic construction of Arithmetic Chern--Simons II, \S3.

The automorphism condition in \eqref{v4-arith:cs:eq:padic-coefficients}
matters. A \(p\)-adic unit exponent acts invertibly on \(\Gamma_1\), but
does not automatically define an automorphism of the full \(\Gamma\). At
\(p=5\), the integer exponent \(2\) is a \(5\)-adic unit and sends \(-1\) to
\(1\).

For an integer exponent \(m\), the full power map is an automorphism
precisely when \(\gcd(m,p(p-1))=1\). One may instead specify a power
automorphism on \(\Gamma_1\) and an independent automorphism on \(\Delta\).
The resulting torsor construction is positive on this domain. It supplies no
complex Hilbert measure without additional scalar and pairing data.

\section{Abelian linking and Wilson correlations}
\label{v4-arith:cs:linking}

The abelian theory has a computable arithmetic correlation, beyond the
definition of a finite sum. Assume now \(\mu_{n^2}\subset F\) and retain the
trivialization \(\zeta\). Put
\[
 V=H^1(X,\mathbb Z/n),\qquad
 V^\vee=\operatorname{Hom}(V,\tfrac1n\mathbb Z/\mathbb Z).
\]
Artin--Verdier duality identifies \(V^\vee\) with
\(\operatorname{Ext}^2_X(\mathbb Z/n,\mathbb G_m)\) and with
\(\operatorname{Cl}(F)/n\). Let \(\delta\) be the Bockstein from
\(0\to\mu_n\to\mu_{n^2}\to\mu_n\to0\). Define \(D:V\to V^\vee\) by
\begin{equation}
 (Dy)(x)=B(x,y)
   :=\operatorname{inv}_\zeta\bigl(x\smile\delta(\zeta_*y)\bigr),
 \label{v4-arith:cs:eq:linking-form}
\end{equation}
where the notation for the invariant includes the indicated coefficient
pairing.

\begin{proposition}[Arithmetic linking pairing]
\label{v4-arith:cs:prop:linking}
The form \(B\) is symmetric. Its radical is \(K=\ker D\). For
\(\xi,\eta\in\operatorname{im}D\), the expression
\[
 \operatorname{ht}_n(\xi,\eta)=B(x,y),
 \qquad Dx=\xi,\quad Dy=\eta,
\]
is well defined and symmetric. A homologically \(n\)-trivial ideal means an
ideal whose class modulo \(n\) lies in \(\operatorname{im}D\). It need not
have zero class modulo \(n\).
\end{proposition}

\begin{proof}
The compatible trivialization modulo \(n^2\) identifies the constant and
roots-of-unity Bocksteins. Their cup-product identity in degree one is
\[
 \delta_2(x\smile\zeta_*y)
   =\delta' x\smile\zeta_*y-x\smile\delta(\zeta_*y).
\]

In this arithmetic setting \(\delta_2=0\). Indeed the \(n\)-adic duality
group \(H^3(X,\mathbb Z_n(1))\cong\mathbb Z_n\) is torsion free, so
reduction in degree two from \(\mathbb Z_n(1)\) to \(\mu_n\) is surjective.
It factors through \(\mu_{n^2}\), which makes the Bockstein zero. Graded
commutativity between degrees two and one now gives \(B(x,y)=B(y,x)\). This
is the argument of Chung--Kim--Kim--Pappas--Park--Yoo, Lemmas~2.1--2.2.

Perfect duality shows that \(B(x,y)=0\) for all \(x\) exactly when \(Dy=0\).
Symmetry gives the same radical in the other variable. Changing either lift
by an element of \(K\) therefore leaves the height unchanged.
\end{proof}

Let \(G=\mathbb Z/n\), and let \(u\in H^1(G,\mathbb Z/n)\) be the identity
character. For the three-class \(c=u\smile\delta'u\), naturality of the
Bockstein identifies the arithmetic action with
\(\operatorname{CS}_c(x)=B(x,x)\). An ideal class \(\xi\in V^\vee\) supplies
the Wilson character
\[
 W_\xi(x)=\exp(2\pi i\langle x,\xi\rangle).
\]
Under class field theory this evaluates the unramified character \(x\) on
the Artin class of the ideal. Its correlation in the quadratic action is
therefore an actual arithmetic finite sum.

\begin{theorem}[Finite Gaussian correlation]
\label{v4-arith:cs:thm:gaussian}
Let \(p\) be odd and let \(V\) be an \(a\)-dimensional \(\mathbb
F_p\)-vector space. Let \(B:V\times V\to\mathbb F_p\) be symmetric with
radical \(K\) of dimension \(b\). Write \(\psi(t)=e^{2\pi it/p}\),
\(r=a-b\), and \(G_p=\sum_{t\in\mathbb F_p}\psi(t^2)\). For \(\lambda\in
V^*\), put
\[
 S(\lambda)=\sum_{x\in V}\psi\bigl(B(x,x)+\lambda(x)\bigr).
\]

If \(\lambda|_K\ne0\), then \(S(\lambda)=0\). Otherwise choose \(z\) with
\(\lambda(x)=B(x,z)\). If \(\overline B\) is the induced nondegenerate form
on \(V/K\), then
\begin{equation}
 S(\lambda)=p^b
       \left(\frac{\det\overline B}{p}\right)
       G_p^r\,\psi\bigl(-B(z,z)/4\bigr).
 \label{v4-arith:cs:eq:gaussian}
\end{equation}

The determinant is taken in a basis of \(V/K\). Its Legendre symbol is
independent of that basis. The empty determinant and \(G_p^0\) are one.
\end{theorem}

\begin{proof}
On a coset of \(K\), the quadratic term is constant. The sum of the additive
character \(\psi\lambda\) over \(K\) is zero unless \(\lambda|_K=0\). In the
latter case, nondegeneracy on \(V/K\) supplies \(z\), unique modulo \(K\).

Since \(2\) is invertible,
\[
 B(x,x)+B(x,z)=B(x+z/2,x+z/2)-B(z,z)/4.
\]
Translation leaves the sum unchanged.

Diagonalize \(\overline B\) over \(\mathbb F_p\). Each diagonal coefficient
\(d\ne0\) contributes \(\sum_t\psi(dt^2)=(d/p)G_p\): counting square roots
of an element gives \(1+(u/p)\), and the substitution \(u\mapsto du\) proves
this identity. The radical contributes \(p^b\). Multiplying proves
\eqref{v4-arith:cs:eq:gaussian}.
\end{proof}

The classical Gauss evaluation is \(G_p=\epsilon_p\sqrt p\), with
\(\epsilon_p=1\) for \(p\equiv1\pmod4\) and \(\epsilon_p=i\) for
\(p\equiv3\pmod4\), for the displayed additive character and the positive
real square root. Thus the magnitude in \eqref{v4-arith:cs:eq:gaussian} is
\(p^{(a+b)/2}\). The formula with this constant is Theorem~1.1 of the
abelian arithmetic Chern--Simons paper.

Apply the theorem to \eqref{v4-arith:cs:eq:linking-form}, after identifying
\(\tfrac1p\mathbb Z/\mathbb Z\) with \(\mathbb F_p\). If each
\(\xi_j=Dz_j\), then \(z=\sum_jz_j\) represents \(\lambda=\sum_j\xi_j\), and
\begin{equation}
 \frac{S(\sum_j\xi_j)}{S(0)}
 =\psi\left(-\frac14\sum_{i,j}
                  \operatorname{ht}_p(\xi_i,\xi_j)\right).
 \label{v4-arith:cs:eq:wilson-linking}
\end{equation}

Here \(S(0)\ne0\). In the action groupoid every abelian field has
automorphism group \(\mathbb Z/p\), so the absolute amplitude is
\(S(\lambda)/p\). This factor cancels in the normalized correlation. The
factor \(1/4\), the radical selection rule and the distinction between
absolute and normalized amplitudes are all forced by the calculation.

At \(p=2\), this completion of the square is unavailable. A quadratic
refinement requires separate treatment.

The linking form also has a class-field description. A \(\mathbb
Z/n\)-torsor on \(X\) determines the invertible \(\mathcal O_F\)-module on
which the group acts through \(\zeta\). Its ideal class is the
class-invariant map \(\Psi:V\to\operatorname{Cl}(F)\). With the Artin
convention of the duality pairing, the map \(D\) is \(\Psi\) followed by
reduction modulo \(n\). Thus the pairing is the evaluation of one unramified
character on the class invariant of the other.

The proof in the same paper, Theorem~4.1, identifies the dual Bockstein with
\(\operatorname{Cl}(F)[n]\to\operatorname{Cl}(F)/n\) by Yoneda products.
This gives the same arithmetic action, with its coefficient trivialization
retained.

\subsection{Local boundary conditions for linking}

There is a relative version on \(U\). For coefficient objects \(G\) on
\(U\), \(F_S\) on its specified local boundary, and a map \(f:F_S\to i^*G\),
its cochain complex is
\[
 C^j(U,G\times_SF_S)
 =C^j(U,G)\oplus C^j(\partial U,F_S)
                    \oplus C^{j-1}(\partial U,i^*G),
\]
with
\[
 d(c,b,a)=(dc,db,da+(-1)^j(fb-\operatorname{res}c)).
\]
For degrees \(i,j\), one cup product is
\begin{equation}
 (c,b,a)\smile(c',b',a')
 =\bigl(c\smile c',b\smile b',
       (-1)^j a\smile f(b')+\operatorname{res}(c)\smile a'\bigr).
 \label{v4-arith:cs:eq:relative-cup}
\end{equation}
Expanding the differential verifies its Leibniz identity. Replacing the
third component by \((-1)^j a\smile\operatorname{res}(c')+f(b)\smile a'\)
changes the product of cocycles by \(d(0,0,(-1)^ia\smile a')\). Thus the two
choices induce the same product in cohomology.

Take \(G=\mathbb Z/n\), \(F_S=\mu_{n^2}\), and let \(f\) be reduction
followed by \(\zeta^{-1}\), with all places dividing \(n\) in \(S\). The
resulting degree-one classes specify characters with chosen local lifts to
\(\mu_{n^2}\). Their connecting map lands in \(H_c^2(U,\mathbb Z/n)\). The
relative duality pairing and this connecting map define the height as above.
The degree-two connecting map vanishes in this setting, and the Bockstein
product identity proves symmetry. This is the relative construction in \S3
of the same paper.

It also gives a local computation without confusing a knot longitude with
inertia. If an ideal class \([I]\) is the connecting image of a character
\(x\), its height against \(J=\prod_{v\notin S}\mathfrak p_v^{e_v}\) is
\[
 \operatorname{ht}_{S,n}(I,J)
       =\sum_{v\notin S}e_v\,x(\operatorname{Art}_v(\varpi_v)).
\]
The Artin maps here use the normalization paired with the chosen duality
invariant. If \(x\) is represented by a Kummer class, each summand is the
corresponding local norm-residue symbol with that same normalization. The
connecting image and reciprocity map must be fixed before replacing this
formula by a Hilbert-symbol convention.

\section{Wilson values and Frobenius}
\label{v4-arith:cs:wilson}

Let \(\sigma:\Gamma\to G\) be a gauge field and let
\(\tau:G\to\operatorname{GL}(V_\tau)\) be a representation over a specified
field. These maps have different sources. For a specified conjugacy class of
loops \(\gamma\) in \(\Gamma\), the Wilson function is
\begin{equation}
 W_{\tau,\gamma}(\sigma)
       =\operatorname{tr}\tau(\sigma(\gamma)).
 \label{v4-arith:cs:eq:wilson-value}
\end{equation}
It is invariant under gauge conjugation.

For a complex gauge representation and an invariant phase \(a\), its
normalized expectation is
\begin{equation}
 \langle W_{\tau,\gamma}\rangle_a
    =\frac{\displaystyle\int_{\mathcal F_G}aW_{\tau,\gamma}}
           {\displaystyle\int_{\mathcal F_G}a},
 \qquad \int_{\mathcal F_G}a\ne0.
 \label{v4-arith:cs:eq:wilson-expectation}
\end{equation}
The denominator condition is part of the definition. A complex phase can
make it vanish: two discrete fields with phases \(1\) and \(-1\) already
give zero.

At trivial action, \eqref{v4-arith:cs:eq:wilson-expectation} is a groupoid
average. It equals a selected field value when the observable is constant on
the support of the weighted sum. Conditioning on one field component also
gives its value. It does not concentrate on a chosen field merely because
the action is zero. On \(\mathcal F_{C_2}(S^1)=C_2//C_2\), each holonomy has
weight \(1/2\), so the sign character has expectation zero and square
expectation one.

For fixed fields, the tensor identity is
\(W_{\tau\otimes\nu,\gamma}=W_{\tau,\gamma}W_{\nu,\gamma}\). Its expectation
is the correlation of the two observables, not generally the product of
their expectations. The sign character example proves this distinction in
dimension one.

\subsection{The unramified quotient}

For a nonarchimedean local field \(F_v\) with residue field \(\mathbb F_q\),
write
\[
 1\longrightarrow I_v\longrightarrow G_{F_v}
       \longrightarrow\widehat{\mathbb Z}\longrightarrow1.
\]
Arithmetic Frobenius is the element of the quotient acting as \(x\mapsto
x^q\) on residue-field extensions. Geometric Frobenius is its inverse. A
lift to \(G_{F_v}\) is well defined modulo inertia. Its image in an
unramified representation is therefore well defined up to conjugacy.

The tame inertia quotient is \(I_v/P_v\cong\prod_{\ell\ne p}\mathbb
Z_\ell(1)\), where \(P_v\) is wild inertia and \(p\) is the residue
characteristic. It is not the unramified quotient \(\widehat{\mathbb Z}\).
In the knot analogy, inertia gives a meridian and a Frobenius lift gives a
longitude. On an unramified rank-\(d\) representation every inertia element
has trace \(d\), while the Frobenius trace can be different.

Nor does \(p\nmid N\) make a representation into
\(\operatorname{GL}_d(\mathbb Z/N)\) unramified. The quadratic extension
\(\mathbb Q_3(\sqrt3)/\mathbb Q_3\) is totally ramified. Its quadratic
character has values \(\{\pm1\}\subset\mathbb F_5^\times\), although
\(3\nmid5\). Unramifiedness is a condition on the kernel of the
representation.

Let \(E_w/F_v\) have residue degree \(f\). Its arithmetic Frobenius maps to
the \(f\)-th power of Frobenius on the unramified quotient of \(F_v\). The
same power rule holds in geometric convention. It gives the corresponding
trace-of-power identity on an unramified carrier. A map between arbitrary
formal neighborhoods at different rational primes has no such
interpretation. Indeed, when \(p\ne q\), there is no unital \(\mathbb
Z\)-algebra map \(\mathbb Z_p\to\mathbb Z_q\): the integer \(q\) is a unit
in the source and not in the target.

The closed immersion \(\operatorname{Spec}\mathbb
F_p\hookrightarrow\operatorname{Spec}\mathbb Z\) has scheme codimension one.
The knot analogy is a comparison of cohomology, fundamental groups and local
duality. It does not change that codimension or furnish a topological
tubular neighborhood with all its gauge-theoretic operations.

\subsection{Coefficient fields and congruences}

The tautological representation of \(\operatorname{GL}_d(\mathbb Z/N)\) has
trace in \(\mathbb Z/N\). For \(N>1\), there is no unital ring map from this
ring to \(\mathbb C\). Thus it is not itself a complex Wilson label. A
complex representation of the same finite group has its own character, which
must be compared to the modular matrix trace by additional data.

A continuous \(\ell\)-adic representation on a finite-dimensional
\(E\)-vector space, with \(E/\mathbb Q_\ell\) finite, admits a stable
\(\mathcal O_E\)-lattice when its image is compact. Fix such a lattice \(L\)
and a uniformizer \(\varpi\). The traces on \(L/\varpi^rL\) are compatible
and satisfy
\[
 \operatorname{tr}\rho(g)
       =\varprojlim_r\operatorname{tr}
              \bigl(\rho(g)\mid L/\varpi^rL\bigr)
       \quad\text{in }\mathcal O_E.
\]
This is one \(\ell\)-adic congruence tower. It does not produce reductions
modulo every integer or complex characters of every congruence group.

If a trace lies in a number field, one can specify its \(\ell\)-adic and
complex embeddings. The topology still matters. Under any field embedding
fixing \(\mathbb Q\), the sequence \(\ell^r\), which tends to zero
\(\ell\)-adically, maps to the same real integers and diverges in \(\mathbb
C\). An abstract scalar embedding therefore cannot turn \(\ell\)-adic
convergence into complex convergence.

An arithmetic Wilson identity now has a precise form. Choose a gauge
representation \(\tau_N\), a Frobenius longitude and a scalar realization.
Then compare its normalized expectation with the trace of the arithmetic
representation. A profinite version additionally needs the actual transition
maps and convergence in the chosen scalar topology. Equality of ranks
supplies none of these identities.

\subsection{Power traces and Euler factors}

Let \(F\) be an endomorphism of a \(d\)-dimensional vector space over a
characteristic-zero field. In formal power series,
\begin{equation}
 \det(1-TF)^{-1}
   =\exp\left(\sum_{m\ge1}\frac{\operatorname{tr}(F^m)}mT^m\right).
 \label{v4-arith:cs:eq:euler-traces}
\end{equation}
After scalar extension, put \(F\) in triangular form. Multiplying
\((1-T\alpha_i)^{-1}\) and expanding its logarithm gives the right side.
Equivalently, differentiating logarithmically and fixing the constant term
proves the identity over the original field.

Thus all power traces determine the Euler polynomial and the semisimple
conjugacy class over an algebraic closure. A single trace does not. The
finite-order matrices \(\operatorname{diag}(1,-1)\) and
\(\operatorname{diag}(i,-i)\) both have trace zero, but their polynomials
are \(1-T^2\) and \(1+T^2\). Even all power traces do not detect a Jordan
extension: \(I_2\) and a nontrivial unipotent Jordan block have trace two in
every positive power.

Characteristic zero is also necessary for this reconstruction from
unrestricted power traces. Over \(\mathbb F_3\), the matrices \(I_3\) and
\(-I_3\) have all power traces zero, while their Euler polynomials are
\(1-T^3\) and \(1+T^3\). Division by \(m\) in
\eqref{v4-arith:cs:eq:euler-traces} cannot be performed in that coefficient
ring. Integral Witt coordinates retain additional data beyond their ghost
traces.

For an unramified \(\ell\)-adic representation at \(v\), \(\ell\ne p\), the
chosen geometric Frobenius gives
\[
 L_v(s,\rho)=\det(1-q^{-s}\rho(\operatorname{Frob}^{\mathrm{geom}}_v))^{-1}.
\]
A Wilson formula for every power would determine this factor after the
coefficient comparison. A Wilson formula for the first trace alone would
not.

For a Weil--Deligne representation \((r,N)\) at a ramified place, the local
factor instead uses the Frobenius-stable space \((\ker N)^{I_v}\):
\[
 L_v(s,r,N)=\det(1-q^{-s}r(\operatorname{Frob}^{\mathrm{geom}}_v)
                                 \mid(\ker N)^{I_v})^{-1}.
\]
The conductor and epsilon factor are further local data. In particular an
epsilon factor also depends on the additive character and measure
convention. It is not inserted into this determinant as an unspecified
correction to the first trace.

If a realization is pure of motivic weight \(m\), its complex Frobenius
eigenvalues have absolute value \(q^{m/2}\). Unitary normalization
multiplies them by \(q^{-m/2}\). This normalization is separate from a
functional-equation center of a global \(L\)-function and from an energy
grading on a Fock space. For example \(\mathbb Q_\ell(-1)\) has motivic
weight two in the cohomological geometric convention. A formal scaling
parameter labeled one does not change it.

\subsection{An elliptic curve at five}

The curve \(E:y^2=x^3-x\) over \(\mathbb F_5\) has affine fiber sizes
\(1,1,2,2,1\) over \(x=0,1,2,3,4\). The point at infinity gives
\(\#E(\mathbb F_5)=8\), so \(a_5=5+1-8=-2\). For \(\ell\ne5\), arithmetic
Frobenius on \(V_\ell E=T_\ell E\otimes\mathbb Q_\ell\) has trace \(-2\) and
determinant \(5\). Its inverse has trace \(-2/5\) and determinant \(1/5\).

The cohomological carrier
\[
 H^1_{\mathrm{\acute et}}(E_{\overline{\mathbb F}_5},\mathbb Q_\ell)
       =\operatorname{Hom}(V_\ell E,\mathbb Q_\ell)
\]
therefore has geometric Frobenius trace \(-2\) and determinant \(5\). Its
local Euler polynomial is \(1+2T+5T^2\). Inertia on the good-reduction
representation is trivial and has trace two, so the meridian and longitude
give different observables in this example.

The Tate-module polarization takes values in \(\mathbb Q_\ell(1)\).
Dualizing gives a cohomological polarization with paired geometric
eigenvalue product \(5\).

For a genus-\(g\) Jacobian the corresponding cohomological rank is \(2g\).
In a Siegel setting the standard and spin representations instead have ranks
\(2g+1\) and \(2^g\). Their Wilson labels and Euler factors must use the
actual representation.

The substitution \(g=0\) gives a rank-zero Jacobian carrier, whose trace is
zero and whose symmetric-algebra character is one. A nontrivial rank-two
arithmetic gauge sector requires a separate representation.

For any finite-dimensional carrier \(V\), the graded trace on
\(\operatorname{Sym}(V)\) is the inverse determinant in
\eqref{v4-arith:cs:eq:euler-traces}. A Heisenberg Fock space with creation
modes at every positive energy has the different character
\(\prod_{r\ge1}\det(1-q^rF)^{-1}\), when the trace is defined. The
oscillator grading must therefore be specified before either expression is
called a boundary character. The arithmetic equality between such a
character and an Euler factor remains a map and trace comparison on the
chosen realization.

\section{Congruence towers and Hilbert limits}
\label{v4-arith:cs:towers}

For a fixed rank \(d\), reduction gives surjections
\(G_N=\operatorname{GL}_d(\mathbb Z/N)\to G_M\) when \(M\mid N\). Their
group limit is \(\operatorname{GL}_d(\widehat{\mathbb Z})\). Passing from
groups to fields, actions and state spaces requires more maps. In
particular, a field at level \(M\) can fail to lift to level \(N\).

For a prime-power step, put \(G_r=\operatorname{GL}_d(\mathbb Z/\ell^r)\),
\(r\ge1\). Then
\begin{equation}
 |G_r|=\ell^{d^2(r-1)}
                  \prod_{j=0}^{d-1}(\ell^d-\ell^j),
 \qquad
 \ker(G_{r+1}\to G_r)=1+\ell^rM_d(\mathbb F_\ell).
 \label{v4-arith:cs:eq:congruence}
\end{equation}
The kernel has the additive multiplication law of \(M_d(\mathbb F_\ell)\),
since \(2r\ge r+1\). It is abelian but generally not central: conjugation by
the permutation matrix exchanging two basis vectors sends \(E_{12}\) to
\(E_{21}\). The cardinality formula follows by counting ordered independent
columns modulo \(\ell\) and then lifting each matrix entry. An invertible
matrix modulo \(\ell\) remains invertible over \(\mathbb Z/\ell^r\).

\subsection{The lifting measure}

Consider an extension of finite groups \(1\to K\to E\xrightarrow{p}G\to1\),
and a finitely presented group \(\Gamma\), or a profinite group with the
specified finite representation sets. For \(\rho:\Gamma\to G\), let
\(\mathcal L_\rho\) be the homotopy fiber of
\(\operatorname{Hom}(\Gamma,E)//E\to \operatorname{Hom}(\Gamma,G)//G\). The
isomorphism to \(\rho\) is part of this fiber.

\begin{proposition}[Lifting mass]
\label{v4-arith:cs:prop:lifting}
Suppose \(K\) is abelian. If \(\rho\) has a lift, conjugation through a lift
defines a \(\Gamma\)-module \(K_\rho\), and
\begin{equation}
 |\mathcal L_\rho|
       =\frac{|H^1(\Gamma,K_\rho)|}{|H^0(\Gamma,K_\rho)|}.
 \label{v4-arith:cs:eq:lifting-mass}
\end{equation}
If there is no lift, the mass is zero. For compatible action phases
\(a_E=p^*a_G\), finite integration gives
\begin{equation}
 Z_E=\sum_{[\rho]}
      \frac{a_G(\rho)|\mathcal L_\rho|}
           {|C_G(\rho(\Gamma))|}.
 \label{v4-arith:cs:eq:tower-fubini}
\end{equation}
\end{proposition}

\begin{proof}
Choose a lift \(\widetilde\rho\). Any other lift has the form
\(z(g)\widetilde\rho(g)\), and the homomorphism law is precisely the
one-cocycle equation for \(z\) with values in \(K_\rho\). The homotopy fiber
is equivalent to these lifts modulo \(K\)-conjugation. Lift an isomorphism
to the target field into \(E\) to make it the identity.

Conjugation changes \(z\) by a one-coboundary. Its stabilizer is
\(H^0(\Gamma,K_\rho)\). This proves the mass formula.

The empty case is immediate. Equation \eqref{v4-arith:cs:eq:tower-fubini} is
Theorem~\ref{v4-arith:cs:thm:fubini} with the phase constant on each fiber.
\end{proof}

For a central kernel \(K=(\mathbb F_\ell)^e\) with trivial action, the mass
is \(|K|^{b_1-b_0}\), where the \(b_i\) are the dimensions of
\(H^i(\Gamma,\mathbb F_\ell)\). For the point it is \(1/|K|\), not one.

Centrality also does not remove the obstruction to lifting. Take \(E=C_4\),
\(G=C_2\), and \(\Gamma=C_2\). The trivial homomorphism to \(G\) has two
lifts, each with two automorphisms in the fiber, so its mass is one. The
identity homomorphism has no lift: an odd element of \(C_4\) cannot be the
image of an element of order two. Its mass is zero. No constant fiber factor
works for both fields.

If all masses equal a specified \(m_{N,M}\), then \(Z_N=m_{N,M}Z_M\). For
the chosen normalization \(\widetilde Z_N=|G_N|^{-\chi}Z_N\), one obtains
\[
 \widetilde Z_N=|K|^{-\chi}m_{N,M}\widetilde Z_M.
\]
The exponent \(\chi\) must come from a specified cohomology theory.
Compatibility requires the displayed factor to be one. The equation uses no
division by a possibly zero partition value, and a scalar equality alone
does not define transition maps on state spaces.

The phases also require compatible levels. Pullback of a fixed cocycle from
a quotient gives one such datum on all finer quotients. If one instead
requires \(c_N=p_{N,M}^*c_M\) for every \(M\mid N\), including \(M=1\),
every normalized three-class is zero, since the bottom group is trivial.
Nontrivial continuous finite-coefficient classes are naturally represented
on a sufficiently fine finite quotient and pulled back above it. More
general transition laws require their own cochains and coherences. A complex
scalar called a level is not by itself such a family of finite-group
three-classes.

\subsection{A positive profinite Hilbert construction}

There is a direct limit of finite Hilbert spaces when the measure and
transition maps are known. Let \(G_{r+1}\twoheadrightarrow G_r\) be a
sequence of finite groups, with kernels \(K_r\), and use uniform probability
measure on each group. Set
\[
 H_r=L^2(G_r,|G_r|^{-1}\text{ counting measure}),\qquad
 (i_rf)(g)=f(p_rg).
\]
Then \(i_r\) is an isometry. Its adjoint is conditional averaging,
\begin{equation}
 (P_rf)(h)=\frac1{|K_r|}\sum_{p_rg=h}f(g),
 \qquad P_ri_r=1.
 \label{v4-arith:cs:eq:conditional-average}
\end{equation}
Both assertions follow by partitioning \(G_{r+1}\) into its equally sized
fibers.

\begin{theorem}[Haar realization]
\label{v4-arith:cs:thm:haar}
For \(G=\varprojlim_rG_r\), normalized Haar measure \(\mu\) has the given
uniform measures as its finite pushforwards. Pullback identifies the Hilbert
completion of \(\operatorname{colim}_r(H_r,i_r)\) with \(L^2(G,\mu)\). The
finite conjugation-invariant subspaces have dense union in the
conjugation-invariant subspace of \(L^2(G,\mu)\).
\end{theorem}

\begin{proof}
The compatible uniform measures assign mass \(1/|G_r|\) to each cylinder
over a point of \(G_r\). They define the normalized invariant measure on the
compact profinite group. Every finite pullback is isometric by this cylinder
formula. Locally constant cylinder functions separate points and are dense
in continuous functions, hence in \(L^2(G,\mu)\). This proves the completion
statement.

The conditional expectation onto level \(r\) commutes with conjugation,
because its fibers and Haar measure do. These orthogonal projections
converge strongly to the identity on the dense union just described.
Applying them to an invariant function proves density in the invariant
subspace.
\end{proof}

This construction uses probability on the group elements. It does not
replace the groupoid measure on arithmetic representations or the action
phase in \eqref{v4-arith:cs:eq:tower-fubini}. To apply it to arithmetic
states one must construct equally precise measures, isometries and adjoints
there.

\subsection{The bounded part of an inverse system}

An inverse diagram always defines a pro-object in its specified category. A
topological inverse limit need not have a Hilbert norm. For example, the
coordinate projections \(\mathbb C^{r+1}\to\mathbb C^r\) have inverse limit
\(\mathbb C^{\mathbb N}\) with its product topology. Every basic
neighborhood constrains finitely many coordinates and contains a whole line
in an unconstrained coordinate. It cannot be a bounded neighborhood for any
norm inducing that topology. Passing to the associated condensed vector
spaces does not supply a missing Hilbert topology.

The orthogonal projections above do have a distinguished Hilbert part. The
following statement also applies to finite-dimensional state spaces once
isometric transitions have been constructed.

\begin{proposition}[Bounded compatible vectors]
\label{v4-arith:cs:prop:martingale}
Let \(H_r\) be Hilbert spaces with isometric maps \(i_r\), let
\(P_r=i_r^*\), and let \(H\) be the Hilbert completion of their increasing
union. Then
\begin{equation}
 H\xrightarrow{\sim}
 \left\{(x_r):P_rx_{r+1}=x_r,
                   \sup_r\|x_r\|<\infty\right\},
 \qquad
 \|x\|^2=\lim_r\|x_r\|^2,
 \label{v4-arith:cs:eq:bounded-limit}
\end{equation}
is an isometric isomorphism, using the projections from \(H\) on the left
and the displayed norm on the right.
\end{proposition}

\begin{proof}
The projections onto the increasing closed subspaces \(H_r\) converge
strongly to the identity of \(H\). They give a compatible sequence with the
asserted norm.

Conversely, view a compatible sequence in the increasing union. For \(s\ge
r\), orthogonality gives
\[
 \|x_s-x_r\|^2=\|x_s\|^2-\|x_r\|^2.
\]
The squared norms increase and are bounded, so this sequence is Cauchy. Its
limit has projections \(x_r\), and its norm is the stated limit. This is
inverse to the first map.
\end{proof}

For the coordinate system, the bounded part is \(\ell^2\), a proper subspace
of the product limit. Specifying this norm condition is a substantive part
of the construction. A Hilbert limit does not ensure trace continuity. The
identity on \(\mathbb C^r\) has trace \(r\), while the identity on
\(\ell^2\) is not trace class. Coherent finite operators alone do not define
an infinite partition trace.

\section{Twisted traces and a finite modular example}
\label{v4-arith:cs:twisted-trace}

The degree-two transgression on circle fields has a categorical recipient.
Let \(\mathcal X\) be a finite groupoid and let \(\tau\) be a normalized
complex multiplicative two-cocycle on its composable arrows. Its twisted
groupoid algebra has basis \(u_a\), one element for each arrow, and
multiplication
\[
 u_a u_b=\tau(a,b)u_{ab}
       \quad\text{when }a,b\text{ are composable},
 \qquad u_a u_b=0\quad\text{otherwise}.
\]
The cocycle identity gives associativity, and the identity arrows give its
unit. A finite skeleton may be used if the groupoid has more than finitely
many objects.

The \emph{inertia groupoid} \(I\mathcal X\) has objects \((x,g)\), with
\(g\in\operatorname{Aut}(x)\), and arrows \(h:(x,g)\to(y,hgh^{-1})\).
Conjugation in the twisted algebra gives a line on this groupoid. On a
one-object component its transport is
\begin{equation}
 u_g\longmapsto
 \frac{\tau(h,g)\tau(hg,h^{-1})}{\tau(h,h^{-1})}
                 u_{hgh^{-1}}.
 \label{v4-arith:cs:eq:inertia-line}
\end{equation}
The same formula applies to composable arrows. Denote this line by
\(L_\tau\). Its degree-zero homology means the direct sum of its fibers
modulo these transport relations.

\begin{proposition}[Finite twisted Hochschild trace]
\label{v4-arith:cs:prop:twisted-hh}
Over \(\mathbb C\), the twisted groupoid algebra satisfies
\begin{equation}
 \operatorname{HH}_0(\mathbb C_\tau[\mathcal X])
       \cong H_0(I\mathcal X,L_\tau),
 \qquad
 \operatorname{HH}_j(\mathbb C_\tau[\mathcal X])=0
                  \quad(j>0).
 \label{v4-arith:cs:eq:twisted-hh}
\end{equation}
The same degree-zero trace is obtained from its category of perfect modules.
\end{proposition}

\begin{proof}
Degree-zero Hochschild homology is the quotient of the algebra by
commutators. Commuting with the identity idempotents kills every non-loop
arrow. For a loop, \([u_h,u_gu_h^{-1}]\) identifies its class with the
conjugate loop in \eqref{v4-arith:cs:eq:inertia-line}. Conversely the cyclic
relation between any two composable arrows is one of these transported loop
relations. This gives the first isomorphism.

Each connected component is Morita equivalent to a twisted algebra of a
finite automorphism group \(H\). It is semisimple: average a linear
projection onto a submodule over conjugation by the invertible operators
\(u_h\). Projective scalar factors cancel under conjugation, so the average
is module-linear. Equivalently, the element
\[
 \frac1{|H|}\sum_{h\in H}u_h\otimes u_h^{-1}
\]
splits multiplication as a bimodule map. The algebra is therefore projective
over its enveloping algebra, and positive Hochschild homology vanishes.
Morita invariance proves the remaining assertions.
\end{proof}

The characteristic assumption is essential. Over \(\mathbb F_p\), the
untwisted algebra of \(C_p\) is \(\mathbb F_p[\epsilon]/(\epsilon^p)\),
whose module of differentials, and hence its first Hochschild homology, is
nonzero. A derived trace cannot then be replaced by its degree-zero inertia
quotient.

For the finite topological theory, \(\mathcal X\) can be the groupoid of
circle fields, with \(\tau\) the transgression of its three-cocycle. Its
category of twisted representations and its Hochschild trace have the types
just described. A trace comparison with an arithmetic Hecke category still
requires that arithmetic category, its coefficient realization and the
actual map on cyclic complexes. Equal dimensions would not define such a
map.

\subsection{The double of \(C_2\)}

The simplest nontrivial modular example can be given without an arithmetic
hypothesis. Let \(A=C_2\), written additively, and let \(\widehat A\) be its
two complex characters. Consider finite-dimensional vector spaces graded by
\(A\times\widehat A\). Simple objects are \((a,\chi)\), tensor product adds
\(a\) and multiplies characters, and the associator is the identity. On
simple objects set
\[
 c_{(a,\chi),(b,\psi)}=\chi(b),\qquad
 \theta_{(a,\chi)}=\chi(a).
\]

The bicharacter law gives both hexagon identities. The balancing identity is
\[
 (\chi\psi)(a+b)
       =\chi(a)\psi(b)\chi(b)\psi(a),
\]
so these data give a ribbon category with all simple dimensions one.

If a simple object has trivial monodromy with every simple, varying \(b\) forces
\(\chi=1\) and varying \(\psi\) forces \(a=0\). Thus the braiding is
nondegenerate.

In the order \((0,1),(0,\mathrm{sgn}),(1,1),(1,\mathrm{sgn})\), the
normalized matrices are
\begin{equation}
 S=\frac12
 \begin{pmatrix}
 1&1&1&1\\1&1&-1&-1\\1&-1&1&-1\\1&-1&-1&1
 \end{pmatrix},
 \qquad
 T=\operatorname{diag}(1,1,1,-1).
 \label{v4-arith:cs:eq:c2-modular}
\end{equation}
Character orthogonality gives \(S^2=1\), and direct multiplication gives
\((ST)^3=1\).

These matrices can also be obtained from the untwisted finite gauge states
on a torus. Its four fields are \((a,b)\in A^2\), each with automorphism
group \(A\). Normalize their delta functions by the measure
\eqref{v4-arith:cs:eq:boundary-hilbert} and Fourier transform in the second
holonomy. The resulting basis is indexed by \((a,\chi)\).

Exchanging the two holonomies acts by
\(S_{(a,\chi),(b,\psi)}=\chi(b)\psi(a)/2\), while a Dehn twist acts by
\(\chi(a)\). This constructs the displayed comparison of finite gauge and
modular data. Its dimension is four, agreeing with \(|\mathcal
F_{C_2}(T^3)|=8/2=4\).

A finite semisimple complex category has finite-dimensional
\(\operatorname{HH}_0\), with one basis vector for each simple object. It
cannot be isomorphic as a vector space to a target containing a nonzero
polynomial Fock space with infinitely many independent creation states. This
excludes that literal finite-category trace lift. It leaves open an infinite
category, an appropriate completion, a selected finite trace sector, or a
compatible family of finite categories. Each choice changes the recipient
and requires its own trace map.

\section{Ordered incidence and local current algebras}
\label{v4-arith:cs:incidence}

An order on places gives an algebraic bar construction before it gives an
arithmetic field theory. Fix a finite ordered set \(S=\{s_1<\cdots<s_r\}\).
An interval partition of \(S\) is determined by a subset of the \(r-1\)
possible cuts. Let \(\mathcal I_S\) have these partitions as objects and one
arrow from a partition to each coarsening obtained by deleting cuts. The
one-block partition \(t_S\) is terminal.

Let \(D_S:\mathcal I_S\to\mathrm{Ch}(k)\) be any covariant diagram of chain
complexes. Its simplicial replacement has
\[
 B_n(D_S)=\bigoplus_{P_0\to\cdots\to P_n}D_S(P_0).
\]

The first face applies the first diagram map. The other faces omit an
intermediate object or the last object. Degeneracies insert identities.
Write \(\partial=\sum_i(-1)^id_i\). With internal homological degree \(q\),
the direct-sum totalization has degree \(n+q\) and differential
\(\partial+(-1)^n d_{\mathrm{int}}\).

\begin{proposition}[Contraction of the incidence bar]
\label{v4-arith:cs:prop:incidence}
The augmentation induced by \(P\to t_S\) is a quasi-isomorphism
\begin{equation}
 \operatorname{Tot}_{\oplus}B_\bullet(D_S)
              \longrightarrow D_S(t_S).
 \label{v4-arith:cs:eq:incidence-contraction}
\end{equation}
It holds for unbounded internal complexes. If the indexing category is the
disjoint union of the \(\mathcal I_S\), its homotopy colimit is
\(\bigoplus_SD_S(t_S)\).
\end{proposition}

\begin{proof}
Include \(D_S(t_S)\) as the terminal summand in degree zero. On a simplex of
degree \(n\), append its unique arrow to \(t_S\) and multiply by
\((-1)^{n+1}\). Denote this map by \(h_n\). The faces cancel in pairs,
except the last face, and in degree zero the terminal augmentation.
Consequently \(\partial h+h\partial=1-i\epsilon\).

The internal differential commutes with appending a terminal arrow. Its
totalization signs in adjacent simplicial degrees are opposite, so the same
identity holds for the total differential. Every vector uses finitely many
simplices in the direct sum. No convergence or lower bound on internal
degree is used. The contraction preserves \(S\), proving the final
assertion.
\end{proof}

For example, if all coefficient complexes are \(k\) and the available places
are \(p,q\), the nonempty components \(\{p\},\{q\},\{p,q\}\) give \(k^3\).
The restricted tensor product of the two copies of \(k\), with unit
reference vectors, is \(k\). The incidence category has no arrows inserting
a new place. Such arrows must be supplied before its colimit can express
restricted tensor assembly.

The simplicial replacement also records partitions of one ordered set, not
all collision diagrams on a curve. Its order is the chosen combinatorial
order on places. A comparison with physical time order or with factorization
on collision strata requires maps that preserve those additional operations.

\subsection{An actual tensor construction of vertex algebras}

There is a positive local construction once current algebras have been
supplied. For every place \(p\), choose a finite-dimensional complex Lie
algebra \(\mathfrak g_p\), an invariant symmetric bilinear form and a level
\(k_p\). The affine brackets are
\[
 [x_m,y_n]=[x,y]_{m+n}
       +m\delta_{m+n,0}k_p(x,y).
\]
The induced vacuum module \(V^{k_p}(\mathfrak g_p)\) is a vertex algebra at
every level. For a finite set \(S\), take its tensor product over \(p\in
S\). Inserting the chosen vacua gives maps for inclusions of finite sets,
and hence a restricted tensor vertex algebra
\begin{equation}
 V_{\mathrm{res}}=\operatorname{colim}_{S\text{ finite}}
                    \bigotimes_{p\in S}V^{k_p}(\mathfrak g_p).
 \label{v4-arith:cs:eq:restricted-current}
\end{equation}

Every vector and every pair of inputs belong to one finite tensor stage.
Their vertex operation, translation and vacuum identities are defined there,
are preserved by enlargement, and therefore define the colimit. Locality for
a pair is the finite tensor locality bound at that stage. This proves the
construction without forming an infinite product of fields on one input.

A conformal vector needs more. For a simple Lie algebra with the standard
normalization, Sugawara requires \(k_p+h_p^\vee\ne0\). A relation between
two shifted levels, \(\kappa_p+\kappa_p^!=0\), does not imply either is
nonzero. Moreover the sum of conformal vectors over infinitely many
nontrivial places need not belong to the algebraic restricted tensor
product. A conformal completion must be specified if that sum is intended.

Satake eigenvalues do not themselves supply these Lie algebras, levels,
vacuum modules or conformal vectors. They can give a finite-dimensional
Frobenius carrier and its symmetric-algebra trace. Identifying that carrier
with a current-algebra sector, and identifying
\eqref{v4-arith:cs:eq:restricted-current} with an arithmetic Hochschild
trace realization, are separate construction problems.

\section{An explicit physical boundary model}
\label{v4-arith:cs:physical}

A physical boundary comparison should identify cochains and operations
before taking a character. The free abelian holomorphic--topological theory
gives such an example. Its base is the actual half-space
\(M=[0,\infty)_t\times\mathbb C_z\), with incoming orientation \(-dt\wedge
dx\wedge dy\), where \(z=x+iy\).

Partial forms use \(dt,d\bar z\), with
\[
 D=dt\,\partial_t+d\bar z\,\partial_{\bar z},\qquad
 \partial=\partial_z,
 \qquad
 \tau_+(a)=\frac1{2\pi i}\int_M^{dt\,dx\,dy}dz\wedge a.
\]
The subscript records the auxiliary positive orientation used in the
integral. Take superfields \(\Phi,\eta\) of total degrees zero and one.

The incoming action and boundary condition are
\begin{equation}
 S_\kappa=-\tau_+\left(\eta D\Phi+
                            \frac\kappa2\eta\partial\eta\right),
 \qquad i^*\eta=0.
 \label{v4-arith:cs:eq:physical-action}
\end{equation}
We fix \(\kappa\in\mathbb C^\times\) and \(\hbar=1\) for the Wilson-jet
comparison below. The action is quadratic. No interacting arithmetic field
theory is being inferred from it.

\subsection{Regular observables and the ending line}

The tangent fields are \(E=\Omega_{\mathrm{HT}}^\bullet e_\Phi
\oplus\Omega_{\mathrm{HT}}^\bullet e_\eta\), with \(|e_\Phi|=0\),
\(|e_\eta|=-1\), and
\begin{align*}
 Q(\alpha e_\Phi)&=D\alpha e_\Phi,\\
 Q(\beta e_\eta)&=D\beta e_\eta+
                  (-1)^{\deg\beta+1}\kappa\partial\beta e_\Phi.
\end{align*}
Compactly supported fields may meet \(t=0\), with the relative condition on
every tangential \(e_\eta\) component. The incoming pairing is
\[
 B_-(\alpha e_\Phi,\beta e_\eta)=-\tau_+(\alpha\beta),
 \qquad
 B_-(\beta e_\eta,\alpha e_\Phi)
                    =(-1)^{\deg\alpha}\tau_+(\beta\alpha),
\]
and vanishes on equal summands. On the suspension, \(|s\xi|=|\xi|-1\), put
\(d(s\xi)=-sQ\xi\) and
\(b_-(s\xi,s\upsilon)=(-1)^{|\xi|+1}B_-(\xi,\upsilon)\).

For \(U=[0,\infty)\times V\), let \(\mathcal
O(U)=\bigoplus_{m\ge0}\operatorname{Sym}_b^m(E_c(U)[1])\). At each field
number, complete the smooth joint kernels in the projective topology at
fixed compact support. Then take the limit over compact supports. The sum
over field numbers is algebraic.

Let \(\Delta_{b_-}\) contract unordered pairs once, with Koszul signs. The
quantum differential is \(d+\Delta_{b_-}\). Smoothness makes its diagonal
contractions defined. A point-supported composite is outside this space.

Write a test field as \(\gamma(t)+dt\,\beta(t)\), and let a superscript
\((0)\) mean Dolbeault degree zero. For \(n\ge0\), the ray and endpoint
derivations are
\begin{align}
 \ell_n(s\xi)&=\frac1{n!}\int_0^\infty
                 \partial_z^n\beta_\eta^{(0)}(t,0)\,dt,
       &|\ell_n|&=1,\notag\\
 r_n(s\xi)&=\frac1{n!}\int_0^\infty
                 \partial_z^n\beta_\Phi^{(0)}(t,0)\,dt,
       &|r_n|&=0,\notag\\
 e_n(s\xi)&=\frac1{n!}\partial_z^n
                    \gamma_\Phi^{(0)}(0,0),
       &|e_n|&=1.
 \label{v4-arith:cs:eq:ray-derivations}
\end{align}
Extend them as constant derivations on kernels.

Integration of the normal derivative gives
\begin{equation}
 [d,\ell_n]=[d,e_n]=0,
 \qquad [d,r_n]=\kappa(n+1)\ell_{n+1}-e_n.
 \label{v4-arith:cs:eq:endpoint-identity}
\end{equation}
The first two identities use the relative boundary condition. The minus sign
in the last uses \(\int_0^\infty\partial_t\gamma(t)\,dt=-\gamma(0)\). All
constant derivations graded-commute with one another and with
\(\Delta_{b_-}\).

For a finite \(F\subset\mathbb Z_{>0}\), introduce the commutative
differential graded coefficient algebra
\begin{equation}
 Y_F=\mathbb C[B_0,B_n:n\in F]
             \otimes\Lambda(c_{n-1}:n\in F),
 \quad
 Q_YB_n=\kappa n c_{n-1},\quad Q_YB_0=Q_Yc_j=0.
 \label{v4-arith:cs:eq:line-coefficients}
\end{equation}
The \(B_n\) have degree zero and the \(c_j\) degree one. Coefficients in
operator expressions act by graded left multiplication. The module has the
regular right \(Y_F\)-action. All tensors with these coefficients are
algebraic.

\begin{proposition}[The endpoint correction]
\label{v4-arith:cs:prop:source-correction}
Put
\begin{align*}
 D_0&=d+\Delta_{b_-}+Q_Y+B_0\ell_0,\\
 L_F&=\sum_{n\in F}\frac{B_n}{\kappa n}r_{n-1},\\
 D_F&=D_0+\sum_{n\in F}
       \left(B_n\ell_n+c_{n-1}r_{n-1}
                        -\frac{B_n}{\kappa n}e_{n-1}\right).
\end{align*}
Then \(D_F=e^{-L_F}D_0e^{L_F}\), so \(D_F^2=0\). The exponential is finite
on each input and fixes the unit and the right coefficient action. Omitting
the endpoint term gives the nonzero square \(\sum_{n\in F}c_{n-1}e_{n-1}\).
\end{proposition}

\begin{proof}
The coefficient differential of \(B_n\ell_n\) gives \(\kappa n
c_{n-1}\ell_n\). The commutator of \(d\) with \(c_{n-1}r_{n-1}\) gives its
negative plus \(c_{n-1}e_{n-1}\), by
\eqref{v4-arith:cs:eq:endpoint-identity}. The endpoint correction cancels
the latter term.

More precisely, \([D_0,L_F]=D_F-D_0\), and its commutator with \(L_F\) is
zero. The conjugation formula follows. Each \(r_n\) lowers field number and
kills the unit, so its exponential terminates and is right-linear. For
\(F=\{1\}\), a suspended \(e_\Phi\) test function equal to one near the
endpoint detects the omitted-term defect as \(c_0\).
\end{proof}

\subsection{All jets with corrected transition maps}

For \(F\subset F'\), coefficient inclusion alone need not commute with the
source differential. The corrected map is
\[
 j_{F,F'}=e^{-(L_{F'}-L_F)}\iota_{F,F'}.
\]
The \(L_F\) commute, so these maps are injective cochain maps and satisfy
\(j_{F',F''}j_{F,F'}=j_{F,F''}\).

For \(Y_\infty=\operatorname{colim}_FY_F\), there is an explicit isomorphism
\begin{equation}
 \operatorname{colim}_F(\mathcal O(U)\otimes Y_F,D_F)
       \xrightarrow{\sim}(\mathcal O(U)\otimes Y_\infty,D_0),
 \qquad [u]_F\longmapsto e^{L_F}u.
 \label{v4-arith:cs:eq:corrected-limit}
\end{equation}

The inverse takes \(v\) to \([e^{-L_F}v]_F\) at any finite stage containing
its coefficients. Enlarging the stage leaves this class unchanged. Each map
is continuous at fixed field number and support, and preserves products on
tangentially disjoint opens.

This is not the literal all-jets differential on an unchanged smooth
observable. On a small disk, choose a compact test whose ray derivatives all
equal one. Its raw source term contains \(\sum_{n>0}c_{n-1}\), which is not
in the polynomial coefficient algebra.
Formula~\eqref{v4-arith:cs:eq:corrected-limit} instead compares the
representatives \(u-\sum_{n\in F}B_n/(\kappa n)\) at different cutoffs.
Their differentials agree under the corrected maps.

\subsection{The boundary period map}

Normal contraction reduces \(D_0\) to the boundary complex
\begin{equation}
 \mathcal A(V)=
 \left(\operatorname{Sym}_b(\Omega_c^{0,\bullet}(V)[1])
              \otimes Y_\infty,
     d_1+\Delta_\mu+B_0\epsilon_0+Q_Y\right),
 \label{v4-arith:cs:eq:boundary-complex}
\end{equation}
where \(d_1(s\alpha)=-s\bar\partial\alpha\), \(\epsilon_0(s\phi)=\phi(0)\)
on functions, and \(\epsilon_0\) vanishes on top forms. The boundary pairing
is \(\mu(s\phi,s\psi)=-\kappa\tau_{\mathbb C} (\phi\partial\psi)\), for a
function \(\phi\) and a \((0,1)\)-form \(\psi\), with \(\tau_{\mathbb
C}(a)=(2\pi i)^{-1}\int dz\wedge a\). The reversed order is given by graded
symmetry.

Here the normal contraction can be constructed with compact support. Choose
a smooth compact profile \(\rho(t)\) of integral one and set
\(F_\rho(t)=\int_0^t\rho(r)\,dr\). On the complementary \(e_\eta\) boundary
component the linear maps are
\begin{align*}
 P_N(\gamma+dt\,\beta)&=\int_0^\infty\beta_\eta(r)\,dr,\\
 I_Ng&=dt\,\rho g+(F_\rho-1)Q_{\mathrm{rel}}g,\\
 h_N(\gamma+dt\,\beta)&=-\int_t^\infty\beta(r)\,dr
          +(1-F_\rho(t))\int_0^\infty\beta_\eta(r)\,dr.
\end{align*}
In the last term the result lies in the \(e_\eta\) summand, and
\(Q_{\mathrm{rel}}=Q-D\). The fundamental theorem of calculus gives
\(P_NI_N=1\) and \(Qh_N+h_NQ=1-I_NP_N\). The correction in \(h_N\) makes its
relative boundary value zero and keeps its normal support compact.

For a boundary form \(a\) of Dolbeault degree \(j\), set
\(i(sa)=sI_N((-1)^jae_\eta)\). On a test field with Dolbeault degree \(j\),
put \(p(s\xi)=s((-1)^jP_N\xi)\), and set \(h(s\xi)=-s h_N\xi\). Then
\(pi=1\), \(dh+hd=1-ip\), and \(hi=ph=h^2=0\).

Let \(H_2\) be the graded symmetrization of \(h\otimes1+ip\otimes h\), and
put \(P_2=-b_-H_2\). Its Hom differential is \(b_--p^*\mu\). Thus
\(p_0=\operatorname{Sym}(p)e^{\partial_{P_2}}\) is the quantum normal
projection to \eqref{v4-arith:cs:eq:boundary-complex}.

The exponential contracts pairs and is finite on each field-polynomial
input. Conjugating the linear contraction gives a homotopy inverse. On the
kernel summand, the symmetric-algebra homotopy is obtained by dividing the
derivation induced by \(h\) by the positive number of kernel factors. These
formulas give a continuous contraction on the stated smooth joint kernels
and commute with disjoint products.

For a disk \(D_R\) centered at zero, define the exterior germ space
\[
 \mathscr P(D_R)=\operatorname{colim}_{r<R}
  \left\{\sum_{n\ge1}a_nw^{-n}:
                     \sum_n|a_n|r^{-n}<\infty\right\}.
\]
Each stage has the displayed weighted \(\ell^1\) norm. The linking maps for
increasing radii are compact.

For a compact top form \(a\), its Cauchy transform is
\[
 \mathcal C(a)(w)=\tau_z\frac{a(z)}{w-z},
\]
expanded outside its support. Set \(\mathcal C=0\) on shifted functions.

The compact Dolbeault complex maps to this germ space by a continuous
quasi-isomorphism. Explicitly, a germ has representative
\(s(\bar\partial\chi)f\) on a slightly larger annulus, where \(\chi\) is one
near the inner disk. If the exterior germ of a top form is zero, Cauchy
convolution is its compact primitive. Cauchy estimates on two separated
radii make both maps continuous. Applying this contraction in each kernel
variable and averaging over permutations proves the statement for every
completed symmetric power.

Let \(k_1(z)=1/z\) as a locally integrable distribution and
\(k_2=-\partial_zk_1\). Put
\[
 P(sa,sb)=\kappa\tau_z\tau_w
                  \bigl(k_2(z-w)a(z)b(w)\bigr),
 \qquad \gamma(sa)=\tau_z(a(z)/z)
\]
on top forms and zero on other degrees. The distributional Cauchy identity
gives \(d_{\mathrm{Hom}}P=\mu\) and \(d_{\mathrm{Hom}}\gamma=\epsilon_0\).
Consequently \(N_{B_0}=e^{\partial_P+B_0\partial_\gamma}\) conjugates the
boundary differential to \(d_1+Q_Y\). Only finitely many contractions act on
each input, and the distributions are paired with smooth kernels in
independent variables.

The resulting period map on the corrected limit is
\begin{equation}
 \Pi[u]_F=\operatorname{Sym}(\mathcal C)
               N_{B_0}p_0e^{L_F}u.
 \label{v4-arith:cs:eq:physical-period}
\end{equation}
Its target is the algebraic sum of completed symmetric powers of \(\mathscr
P(D_R)\), tensored algebraically with \(Y_\infty\).

\begin{theorem}[Rotation-finite boundary comparison]
\label{v4-arith:cs:thm:physical-period}
Rotate \(z\) and its form factors geometrically, with an additional weight
one on \(e_\Phi\). Give \(B_n\) and \(c_{n-1}\) weight \(n\), and \(B_0\)
weight zero. Let \(\mathcal J_{\mathrm{rot}}(D_R)\) be the algebraic sum of
the circle eigenspaces of the corrected source. Then
\eqref{v4-arith:cs:eq:physical-period} restricts to a continuous right
\(Y_\infty\)-linear quasi-isomorphism
\begin{equation}
 \Pi_{\mathrm{rot}}:\mathcal J_{\mathrm{rot}}(D_R)
       \longrightarrow
 (F_{\mathrm{pol}}\otimes Y_\infty,Q_Y),
 \qquad F_{\mathrm{pol}}=\mathbb C[x_1,x_2,\ldots].
 \label{v4-arith:cs:eq:polynomial-period}
\end{equation}
It sends \(1\) to \(1\), is compatible with centered disk inclusions, and
carries separated contour currents to
\begin{equation}
 J_{-n}=x_n,\qquad J_n=\kappa n\partial_{x_n}\quad(n>0),
 \qquad J_0=B_0.
 \label{v4-arith:cs:eq:current-modes}
\end{equation}
For every algebraic differential graded left \(Y_\infty\)-module \(L\),
ordinary tensoring over \(Y_\infty\) computes the derived tensor product and
preserves this quasi-isomorphism.
\end{theorem}

\begin{proof}
The corrected transition isomorphism, normal contraction, Gaussian
conjugation and compact Dolbeault contraction prove the comparison on the
full germ target. They commute with rotations. A monomial \(w_1^{-n_1}\cdots
w_m^{-n_m}\) has weight \(-\sum n_i\). For a fixed sum, only finitely many
symmetric monomials occur, and \(m\) is bounded by that sum. Thus the
finite-circle-spectrum part of the germ target is exactly
\(F_{\mathrm{pol}}\). Each input has finitely many coefficient monomials, so
adjoining \(Y_\infty\) preserves this conclusion.

To restrict the quasi-isomorphism, take a cocycle of finite circle spectrum
in its acyclic mapping cone. Choose a primitive in the full cone and apply
the finite Fourier projector onto the cocycle's weights. The rotated compact
support stays in one larger compact disk. The integral therefore exists in a
single smooth-kernel or germ stage and commutes with the differential. It
gives a primitive of finite spectrum.

A current smeared on an outer circle has period
\[
 \sum_{n>0}x_nz^{n-1}+B_0z^{-1}
                     +\kappa\sum_{n>0}n z^{-n-1}\partial_{x_n}.
\]
The uncontracted insertion gives the first sum, its charge contraction gives
the middle term, and \(P(sa_z,sa)=-\kappa\partial_z\mathcal C(a)(z)\) gives
the last. Contour coefficients yield \eqref{v4-arith:cs:eq:current-modes}.
Moving the circle without crossing the inner support gives a cochain
homotopy by Stokes' formula.

Finally filter the source by field number. Its differential does not raise
that number. Every associated graded piece is a complex of vector spaces
tensored with the regular coefficient algebra, and is homotopically flat
over that algebra. The exhaustive, bounded-below filtration and exact
filtered unions give homotopical flatness of the source. The target is free
as a graded coefficient module. Tensoring the quasi-isomorphism with any
algebraic differential graded module therefore preserves it, including
unbounded modules.
\end{proof}

\subsection{The Gaussian cap}

On the boundary write \(\Phi_\Sigma=\phi+d\bar z\,\phi_b\) and
\(\eta_\Sigma=c+d\bar z\,a_b\). Here \(c(z)\) is the boundary ghost. The
\(c_j\) in \eqref{v4-arith:cs:eq:line-coefficients} are line coefficients.
Put \(x_n=\partial_z^{n-1}\phi(0)/(n-1)!\). The incoming boundary charge,
including the endpoint, is
\begin{align*}
 S_{\mathrm{in}}^F
  &=-\frac1\pi\int c(\partial_{\bar z}\phi+
                           \kappa\partial_z a_b)\,dx\,dy+\mu_F(0),\\
 \mu_F(0)
  &=B_0c(0)+\sum_{n\in F}
       \left(\frac{B_n\partial_z^nc(0)}{n!}+c_{n-1}x_n\right).
\end{align*}

The coefficient differential \(Q_Y\) is external to the boundary Poisson
bracket, since the line coefficients are Casimirs. Its contribution
\(\pi\kappa n c_{n-1}\partial_z^ng(0)/n!\) to the square on \(\phi_b(g)\)
cancels the contribution \(-\pi\kappa c_{n-1}\partial_z^ng(0)/(n-1)!\) from
the line-ghost term. The remaining square is zero on every field generator.

Let \(\mathcal P_{\mathrm{sm}}(V)\) be the polynomial functionals of
\(\phi,\phi_b\) with compact smooth kernels. For a finite jet set \(F\)
and \(N\ge\max(F\cup\{1\})\), set
\[
 W_{F,N}=Y_F[[x_1,\ldots,x_N]],\qquad
 \mathcal H_{\mathrm{in}}^F
   =\operatorname{colim}_N
       \bigl(\mathcal P_{\mathrm{sm}}(V)\otimes_{\mathbb C}W_{F,N}\bigr).
\]
Completion is by jet weight, \(\operatorname{wt}x_n=n\), before taking the
union in \(N\). Each weight coefficient is polynomial in the line
coefficients. The tensor with smooth functionals is algebraic.

In the incoming polarization, momenta act as
\(\widehat a_b=\pi\delta/\delta\phi\) and
\(\widehat c=\pi\delta/\delta\phi_b\), with derivatives on the right.
Quantizing the charge defines \(\Omega_{\mathrm{in}}^F\) on this space.
A mixed contraction differentiates one of the finitely many jet variables
and lowers weight by at most \(N\). Its other factor takes finitely many
derivatives of a smooth kernel at the junction. Smooth-field contractions
restrict smooth joint kernels to diagonals. The operator is therefore
defined and continuous at each fixed \(N\).

The square of \(Q_Y+\Omega_{\mathrm{in}}^F\) is zero. Each pair of charge
terms shares at most one position--momentum contraction, so its commutator
quantizes the classical bracket. The two-bulk-term contribution pairs a
symmetric mixed-derivative kernel with two antisymmetric ghosts and
vanishes. The endpoint bracket cancels \(Q_Y\mu_F(0)\) by the displayed
factorial identity.

On the pure junction subspace \(W_F=\operatorname{colim}_NW_{F,N}\), every
momentum term contains \(\delta/\delta\phi_b\) and vanishes. The resulting
differential is
\[
 D_{\mathrm{cap}}=Q_Y+\sum_{n\in F}x_nc_{n-1}.
\]
Its multiplication term is the ghost part of the endpoint charge. Omitting
it would violate the relative master equation of the coupled boundary and
line system.

Consider a cap \(e^f\) with \(f\) polynomial in finitely many \(x_n\),
normalized by \(f|_{x=0}=0\). It is closed precisely when
\[
 f=-\sum_{n\in F}\frac{B_nx_n}{\kappa n}+h,
 \qquad h\in\mathbb C[B_0,x_1,x_2,\ldots].
\]
Indeed \(D_{\mathrm{cap}}e^f=e^f (Q_Yf+\sum x_nc_{n-1})\), so the displayed
equation integrates the coefficient derivatives.

Under rotation, \(x_n\) has weight \(-n\). The first sum has weight zero and
every nonconstant monomial of \(h\) has negative weight. Thus the unique
rotation-invariant normalized cap in this class is
\begin{equation}
 G_F=\exp\left(-\sum_{n\in F}\frac{B_nx_n}{\kappa n}\right).
 \label{v4-arith:cs:eq:cap}
\end{equation}

The constant state is not closed when \(F\ne\varnothing\). Gluing the linear
exponent of \(G_F\) by the BV pairing produces exactly \(-\sum(B_n/(\kappa
n))e_{n-1}\), the endpoint correction of
Proposition~\ref{v4-arith:cs:prop:source-correction}. Changing its sign
leaves square \(2\sum c_{n-1}e_{n-1}\).

For all jets, use the degreewise creation-weight completion
\(\prod_{N\ge0}(F_N\otimes Y_\infty^d)\), where \(F_N\subset
F_{\mathrm{pol}}\) has creation weight \(N\). In it set
\[
 G=\exp\left(-\sum_{n>0}\frac{x_nB_n}{\kappa n}\right),
 \qquad \mathsf P_\kappa=G(F_{\mathrm{pol}}\otimes Y_\infty).
\]
Every weight coefficient is finite. On this specified image, multiplication
by \(G\) is an invertible cochain map from \(Q_Y\) to
\(D_G=Q_Y+\sum_{n>0}x_nc_{n-1}\). Hence \(G\Pi_{\mathrm{rot}}\) sends the
supported unit to the cap and carries the currents to
\[
 \widehat J_{-n}=x_n,\qquad
 \widehat J_n=\kappa n\partial_{x_n}+B_n,\qquad
 \widehat J_0=B_0.
\]

Coefficientwise differentiation proves \(D_G(Gf)=GQ_Yf\) and \(\widehat
J_n(Gf)=G\kappa n\partial_{x_n}f\). The chosen image \(\mathsf P_\kappa\),
the full completion and the raw polynomial complex are different carriers.

The same formulas detect a possible Ward error. For \(m>0\), the raw
operator \(O_m=-(\widehat J_m+B_m)\) has
\[
 [D_G,O_m]=-\kappa m c_{m-1},\qquad
 O_mG=-B_mG.
\]
Adding the coefficient term \(B_m\) gives the cochain map \(-\widehat J_m\).
A residue identity for a raw insertion therefore cannot replace the
correction that makes the operation a cochain map.

The order of state completion also matters. If all jet variables are
completed before adjoining smooth functionals, take a smooth test equal to
\((1-z)^{-1}\) near the endpoint. Applying the two-momentum part of the
boundary operator to \(\phi_b(g)\sum_{n=1}^Mx_n\) has scalar term
\(\pi\kappa M(M+1)/2\). It does not stabilize even in weight zero. The
finite-jet completion used to construct the cap avoids this undefined
operation.

\subsection{Rigidity in the regular polynomial class}

The cap also explains which local choices survive rotation. Let \(\mathfrak
J\) be the complex of constant derivations obtained by taking finitely many
normal, holomorphic and antiholomorphic derivatives of test fields at the
endpoint. Its differential is \([d,-]\). The relative boundary condition
makes its cohomology
\[
 H^1(\mathfrak J)=\bigoplus_{n\ge0}\mathbb C[e_n],
 \qquad H^j(\mathfrak J)=0\quad(j\ne1).
\]

To see this, take the restricted dual of formal field jets, graded by normal
and antiholomorphic order and rotation weight. Each piece is finite. The
relative \(e_\eta\) normal complex \(t\mathbb C[[t]]\xrightarrow{\partial_t}
\mathbb C[[t]]dt\) is acyclic. The unrestricted \(e_\Phi\) normal complex
has constant cohomology, and its formal Dolbeault complex has cohomology
\(\mathbb C[[z]]\). Dualizing gives precisely the holomorphic evaluation
classes \([e_n]\).

Consider junction corrections of degree one in \((Y_F\otimes\mathfrak
J)\oplus Y_F\), the second summand acting by multiplication. Write
\(X_F^0=-\sum(B_n/(\kappa n))e_{n-1}\). The square-zero corrections are
exactly
\begin{equation}
 X=X_F^0+\partial Y+\sum_{n\ge0}h_ne_n-Q_Y(q),
 \label{v4-arith:cs:eq:junction-rigidity}
\end{equation}
where \(\partial=Q_Y+[d,-]\), \(Y\) is a degree-zero coefficient-valued
constant derivation, the sum is finite, \(h_n\in\mathbb C[B_0]\), and
\(q\in\mathbb C[B_0,B_n:n\in F]\). Indeed the square is
\(\partial(X-X_F^0)\). The polynomial de Rham contraction gives
\(H(Y_F)=\mathbb C[B_0]\), and the preceding jet calculation identifies all
its degree-one classes. This proves
\eqref{v4-arith:cs:eq:junction-rigidity}.

The term \(\partial Y\) is removed by the finite conjugation \(e^Y\), which
fixes the unit and the regular right action. The classes \(h_ne_n\) have
weight \(-(n+1)\), and a nonzero \(Q_Y(q)\) has positive coefficient weight.
Rotation invariance therefore forces both to vanish and permits a
weight-zero \(Y\). In this regular polynomial class, a rotation-invariant
junction is isomorphic to the one used in
Theorem~\ref{v4-arith:cs:thm:physical-period}. The conclusion concerns this
local source and its admissible corrections. It is not a generation theorem
for every boundary module.

A coefficient twist can change the polynomial ground. With correction
\(X_F^0-Q_Y(q)\) and \(q\notin\mathbb C[B_0]\), the period differential
becomes \(Q_Y-Q_Y(q)\). A polynomial ground killed by all positive currents
would satisfy \(\partial_{B_n}f=(\partial_{B_n}q)f\).

For a variable with \(\partial_{B_n}q\ne0\), the right side has total
polynomial degree at least that of nonzero \(f\), while the left side has
smaller degree. Thus no normalized polynomial ground exists.

Completing in the positive \(B_n\) permits the ground
\(e^{q-q|_{B_{>0}=0}}\). This repairs the comparison on a different
coefficient space and changes the selected vector.

\section{Sewing states and evaluating the formal comparison}
\label{v4-arith:cs:sewing}

For an algebraic differential graded line module \(L\), write
\(\Phi_\kappa(L)=F_{\mathrm{pol}}\otimes L\) with its differential and the
currents \eqref{v4-arith:cs:eq:current-modes}. Two endpoints carry two
copies of the coefficient algebra. On \(U=M\otimes N'\), their zero modes
commute. Put \(N=B'_0B''_0/\kappa\). The state family for boundary fusion
has carrier \(F_{\mathrm{pol}}\otimes U\) and connection
\begin{equation}
 \nabla_{\mathrm{st}}=d+\rho(N)\,d\log z,
 \label{v4-arith:cs:eq:state-connection}
\end{equation}
where \(\rho\) denotes the coefficient action and \(z\ne0\) is the
separation.

For scalar charges \(a,b\), its normalized current operation on the first
ground state is
\begin{equation}
 K_{a,b}(z)f(x)|b\rangle
 =\exp\left(\frac a\kappa
                 \sum_{n>0}\frac{x_nz^n}{n}\right)
       f(x_n-az^{-n})|a+b\rangle.
 \label{v4-arith:cs:eq:normalized-sewing}
\end{equation}
The output is completed in creation weight. Each weight coefficient on a
polynomial input is a Laurent polynomial in \(z\). Commuting \(\kappa
m\partial_{x_m}\) and \(x_m\) through the two factors verifies \([
J_m,K_{a,b}(z)]=az^mK_{a,b}(z)\), where the commutator uses the input and
output current actions.

Negative-current descendants of the first input are determined by the
normal-ordered insertion rule
\[
 I(J_{-n}u,z)=
 \frac{\partial_z^{n-1}J^-(z)}{(n-1)!}I(u,z)
 +I(u,z)\frac{\partial_z^{n-1}J^+(z)}{(n-1)!}.
\]
Here \(J^-\) contains negative modes and \(J^+\) nonnegative modes.
Truncation on polynomial inputs and completion of output weight make each
expression defined.

Removing \(K_{a,b}\) from an operation satisfying these identities leaves a
map commuting with all nonzero oscillators. It is determined by the image of
the ground state. More generally, for a coefficient module with locally
nilpotent positive modes and positive-mode truncation on each vector, the
oscillator decomposition is \(F_{\mathrm{pol}}\otimes\Omega_+\), where
\(\Omega_+=\bigcap_{m>0}\ker J_m\). Thus normalized Ward operations are
represented by the state family above.

On the charge-\(a\) Fock space, translation is
\[
 L_{-1}^{(a)}=\frac a\kappa x_1+
                    \sum_{n>0}n x_{n+1}\partial_{x_n}.
\]
Substituting \eqref{v4-arith:cs:eq:normalized-sewing} into translation
covariance makes its ground coefficient satisfy \(\partial_z h=(N/z)h\). Its
connection is therefore \(d-\rho(N)d\log z\). Formally at \(z=1\), its
normalized horizontal coefficient is \(\exp(N\log(1+(z-1)))\).

For scalar charges it becomes \(z^{ab/\kappa}\) on a logarithmic cover. A
counterclockwise half-loop transports states by \(e^{-\pi iab/\kappa}\) and
operations by \(e^{\pi iab/\kappa}\). The ordinary tensor product of two
level-\(\kappa\) Heisenberg representations has level \(2\kappa\) under the
diagonal current. The fusion just specified remains at level \(\kappa\) and
adds the charges.

\subsection{The jet twist}

The classical Gauss constraint alone does not impose the holomorphic Ward
identity. On a disk of radius \(R\) with two smaller disks of radii
\(r',r''\) removed, the partial connection \(a=d\bar t\) is closed. With
zero sources its mode-one contour defect is \(\kappa(R^2-r'^2-r''^2)\). This
follows from \(\oint(t-c)d\bar t=-2\pi i r^2\) around a counterclockwise
circle. At \(\kappa=2\), \(R=2\), \(r'=r''=1/4\), the defect is \(31/4\).
The quantum comparison above uses the constructed current operations, not
that constraint alone.

The line-side family has an additional jet twist. In the two coefficient
algebras set
\begin{equation}
 s(z)=\frac1\kappa
   \sum_{p+m\ge1}(-1)^{p+1}\binom{p+m}{p}
       \frac{B'_pB''_m}{(p+m)z^{p+m}},
 \qquad r=Q_Ys,
 \qquad h=\frac Nz+\partial_zs.
 \label{v4-arith:cs:eq:jet-twist}
\end{equation}
Both indices are nonnegative. Coefficients outside a chosen finite jet set
are zero when that restriction is imposed. For all jets the sum is a power
series of positive order in \(q=z^{-1}\).

Each coefficient is a finite sum. Its generating scalar is
\(\kappa^{-1}\log(1+(u-v)/z)\), with \(u^p v^m\) replaced by \(B'_pB''_m\).
In particular \(Q_Ys=r\) and \(\partial_zs=h-N/z\), with all signs fixed.

\begin{theorem}[Formal line-to-state comparison]
\label{v4-arith:cs:thm:formal-sewing}
Replace the coefficient complex \(U\) degreewise by \(U^d((q))\), the
Laurent series with a finite lower bound and arbitrary coefficient vectors.
On \(F_{\mathrm{pol}}\otimes U((q))\), the map
\begin{equation}
 E(q)=e^{-\rho(s(z))}
 \label{v4-arith:cs:eq:sewing-map}
\end{equation}
is an isomorphism from the differential \(d_U\) and connection
\(\partial_z+\rho(N)/z\) to \(d_U+\rho(r)\) and \(\partial_z+\rho(h)\). It
commutes with the currents and their rooted residue operations, and
preserves the identity action of the vertex vacuum. It need not fix a chosen
cyclic vector.
\end{theorem}

\begin{proof}
The differential graded module identity gives \([d_U,\rho(s)]=\rho(r)\). The
even operator \(\rho(s)\) commutes with \(\rho(r)\) and \(\rho(h)\), because
the coefficient algebra is graded commutative. Thus
\[
 e^{\rho(s)}(d_U+\rho(r))e^{-\rho(s)}=d_U,
 \qquad
 e^{\rho(s)}(\partial_z+\rho(h))e^{-\rho(s)}
                         =\partial_z+\rho(N)/z.
\]

There are no further commutators. Positive \(q\)-order makes each
coefficient of the exponential and its inverse finite. This proves the
formulas on the stated Laurent completion for every algebraic module.

The map acts only on coefficients and commutes with all oscillator modes and
the total zero mode. It therefore commutes with normal-ordered current
operations and their residues. Applied to a ground vector \(u\), it gives
\(e^{-\rho(s)}u\), which proves the final distinction.
\end{proof}

For finitely many endpoints, the pair twists commute. On a common formal
separation domain, the exponential of their sum is independent of
parenthesization. After removing it, the state connection is
\[
 d+\sum_{i<j}\frac{B^{(i)}_0B^{(j)}_0}{\kappa}
                                      d\log(z_i-z_j).
\]
This proves formal refinement coherence on that domain. It does not evaluate
an arbitrary all-jets series at numerical separations.

\subsection{An exact evaluation criterion}

Fix finitely many jets and write \(\rho(s)=S(q)=\sum_{k=1}^KA_kq^k\), with
commuting degree-zero operators \(A_k\). In this paragraph, an exponential
is defined when its partial sums on each vector converge within a
finite-dimensional subspace. The operators are \emph{jointly locally finite}
when \(\mathbb C[A_1,\ldots,A_K]u\) is finite dimensional for every \(u\).

\begin{proposition}[Algebraic evaluation]
\label{v4-arith:cs:prop:evaluation}
At a fixed nonzero separation, \(e^{-S(q)}\) is defined on algebraic states
exactly when \(S(q)\) is locally finite. The formal map preserves
\(U\otimes\mathbb C((q))\) exactly when the \(A_k\) are jointly locally
finite. This is equivalent to evaluation being defined at every nonzero
separation. Evaluation at \(K\) distinct nonzero values of \(q\) already
suffices. On this domain a removable singularity of the comparison at
\(z=0\) occurs exactly when every \(A_k\) vanishes, in which case the
comparison is the identity.
\end{proposition}

\begin{proof}
Successive differences of partial sums on \(u\) recover \(S(q)^nu/n!\).
Their finite-dimensional convergence is therefore equivalent to local
finiteness, and the inverse is the matrix exponential on each finite orbit.
These orbits may be enlarged by those of \(d_Uu,\rho(r)u,\rho(r)d_Uu\). The
relations \([d_U,S]=\rho(r)\), \([S,\rho(r)]=0\) and \(\rho(r)^2=0\) then
justify the cochain identity on a finite-dimensional stable space.

A vector-valued series lies in the algebraic scalar extension precisely when
its coefficients span a finite-dimensional space. A finite joint orbit
clearly suffices.

For necessity, identify the cyclic coefficient module of a nonzero \(u\)
with the quotient
\[
 B=\mathbb C[t_1,\ldots,t_K]/\operatorname{Ann}(u).
\]
If the coefficients of \(\exp(-\sum q^kt_k)\) span a finite space,
evaluation at characters of \(B\) puts all scalar functions
\(\exp(-\sum q^k\chi(t_k))\) in a finite-dimensional space.

Exponentials of distinct polynomials with zero constant term are linearly
independent. Choose a ray avoiding the finitely many leading real-part
coincidences. Divide a dependence by its fastest growing term and let the
radius tend to infinity. Proceed inductively.

Thus \(B\) has finitely many characters. The Nullstellensatz makes each
polynomial \(\prod_\chi(t_k-\chi(t_k))\) nilpotent in \(B\). A power is a
monic relation for \(t_k\), so reducing all exponents shows that \(B\) is
finite dimensional.

If \(S(q_j)\) are locally finite at \(K\) distinct nonzero \(q_j\), their
commuting joint orbit is finite on each vector, by their monic relations.
The matrix \((q_j^k)\) is invertible, so each \(A_k\) is a linear
combination of these operators. This proves joint local finiteness. Finally
on each finite joint orbit \(e^{-S(q)}\) is entire in \(q\). Removability at
\(z=0\) makes its entries bounded at infinity. Liouville's theorem makes it
the constant identity, and the formal logarithm gives \(S=0\).
\end{proof}

The coefficientwise completion is strictly larger in a basic example. Let
\(M=\mathbb C[b]\otimes\Lambda(c)\), \(d_M=\kappa c\partial_b\), with
\(B_0=0\), \(B_1=b\), \(c_0=c\), and all higher coefficients zero. Take a
second module of scalar charge one. Then \(s=b/(\kappa z)\), \(r=c/z\), and
\[
 e^{-s}1=\sum_{n\ge0}
              \frac{(-1)^nb^nq^n}{n!\kappa^n}
\]
lies in \(M((q))\) but not in \(M\otimes\mathbb C((q))\).

At fixed \(z\ne0\), the untwisted degree-zero cohomology is \(\mathbb C\).
The twisted equation \(\kappa f'(b)+f(b)/z=0\) has no nonzero polynomial
solution, by its highest coefficient. The two algebraic complexes are not
quasi-isomorphic. The formal theorem remains valid on its larger coefficient
space.

\subsection{Coincidence and nested regions}

Universal coefficients also admit a distributional interpretation with fixed
normalization. For \(k\ge1\), define
\[
 T_k=\frac{(-1)^{k-1}}{(k-1)!}
                   \partial_z^{k-1}(1/z).
\]
It extends \(z^{-k}\) and satisfies \(\bar\partial T_k=
\pi(-1)^{k-1}\partial_z^{k-1}\delta_0/(k-1)!\). At finite jet set, complete
the coefficient algebra in the positive-jet ideal. Each coefficient of
\(e^{-s}\) is a finite linear combination of powers \(z^{-k}\), so this
formula extends it coefficientwise and fixes its contact derivatives. It
does not multiply distributions with the same singular support or evaluate
the whole coefficient series on an arbitrary module.

For a smooth radial approximate identity \(\rho_r\) of mass one, these fixed
coefficients satisfy \(T_k*\rho_r\to T_k\) with error \(O(r^2)\) on each
compact smooth test. Taylor expansion of the test function has no
first-moment contribution by radial symmetry. A finite-order distribution
bound controls the second remainder. This gives a precise local smoothing
statement with the same contact normalization.

Composition of Laurent expansions has its own domain. For example,
\[
 \frac1{(z_1-z_2)^2}
  =\sum_{n\ge0}(n+1)\frac{z_2^n}{z_1^{n+2}}
       \qquad (|z_2|<|z_1|).
\]
The series converges with all derivatives on compact subsets of this nested
region.

No distributional convergence holds on an open set crossing the boundary. In
the ratio variable, take a radial test centered at a point of modulus
greater than one. Holomorphicity makes its integral the value at the center,
which does not tend to zero. The existence of the distribution
\(T_2(z_1-z_2)\) does not enlarge the convergence domain of this particular
composition series. Wick products inherit the corresponding nested domains,
while their formal refinement identities use the common coefficient
completion specified above.

Finally, the comparison map and the binary operation have different
coincidence behavior. All positive coefficients may act by zero, making
\(E=1\), while the current operation still has a pole. For scalar charges
with \(ab/\kappa\notin\mathbb Z\), the state connection has monodromy
\(e^{-2\pi iab/\kappa}\), and \(z^{ab/\kappa}K_{a,b}\) is defined on
branches. It has no single-valued distributional extension across a disk
about zero agreeing with those branches. These are conditions on the
operation and its coefficient local system, not a failure of the formal
coefficient isomorphism.

\section{Arithmetic bulk and boundary comparisons}
\label{v4-arith:cs:comparisons}

The finite arithmetic torsor, the topological state sum and the free
physical boundary complex have now been constructed on their stated domains.
An arithmetic bulk theory relating them needs a common geometric and
categorical recipient. Write \(\mathscr X\) for a specified arithmetic site
and \(j:\mathscr B_\infty\to\mathscr X\) for a proposed archimedean boundary
morphism. An orientation consists of a dualizing coefficient object and a
trace in its specified degree. In a suitable coefficient theory, it is a map
\(R\Gamma_c(\mathscr X,\omega_{\mathscr X})\to k[-3]\). It must be
compatible with the restriction and relative duality used by the action.

At the boundary, distinguish a linear category \(\mathcal B\), its
Hochschild complex, a marked trace class, and a chosen state realization
\(E_\infty\). For example, the intended primary arithmetic datum can be a
marked class in the Hochschild trace of a specified Iwahori--Hecke category.
Extracting \(E_\infty\) from that trace requires its degree, coefficient
theory and completion. The trace class is not a vertex algebra or a category
of its own modules. A boundary comparison must give maps of these actual
objects preserving the marked class and the operations under consideration.

Three arithmetic identities then have different content. The boundary
identity compares the designated trace or state realizations. The Wilson
identity compares normalized field integrals with Frobenius traces in the
chosen coefficients. The partition identity compares a defined scalar trace
with a completed \(L\)-function.

For the Riemann normalization the latter function is
\[
 \xi(s)=\tfrac12s(s-1)\pi^{-s/2}
                       \Gamma(s/2)\zeta(s).
\]
Its polar-removal factor cannot be omitted when the comparison is made with
\(\xi\) rather than \(\pi^{-s/2}\Gamma(s/2)\zeta(s)\). An identity as a
function of \(s\) needs an analytic family of observables or traces and a
common domain. Renaming a fixed finite-group level \(s\) does not produce
such a family.

\subsection{Finite towers}

A tower construction starts with actual finite field functors and coherent
transition maps on their states and actions.
Proposition~\ref{v4-arith:cs:prop:lifting} computes the measure required by
a reduction. A coherent inverse diagram defines a pro-object in the chosen
functor category. If a represented limit exists and its composition maps
preserve the finite gluing identities, it defines the corresponding limiting
assignment.

A Hilbert realization can be constructed under the isometric and boundedness
hypotheses of Theorem~\ref{v4-arith:cs:thm:haar} and
Proposition~\ref{v4-arith:cs:prop:martingale}. Arithmetic field spaces need
their own measures and maps satisfying those hypotheses. The boundary,
Wilson and partition comparisons are further identities on that limit.

Even the untwisted closed arithmetic example distinguishes the finite sum
from a zeta partition formula. The \'{e}tale fundamental group of
\(\operatorname{Spec}\mathbb Z\) is trivial: a nontrivial finite extension
unramified at all finite primes would contradict the discriminant bound. Its
finite \(G\)-field groupoid is \(BG\), so its untwisted mass is \(1/|G|\).
Additional ramification, boundary or analytic data are needed for a
nonconstant arithmetic partition family.

\subsection{Modular categories}

A finite semisimple ribbon category with nondegenerate braiding gives the
algebraic input to a Reshetikhin--Turaev theory, together with its anomaly
and tangential structure. The \(C_2\) example above provides explicit such
data. To apply \(Z_{\mathrm{RT}}\) to arithmetic objects, construct an
arithmetic bordism category. Supply also a symmetric monoidal functor \(C\)
into the appropriate topological bordism category. Their composition
\(Z_{\mathrm{RT}}C\) is then an arithmetic assignment.

This composition preserves the monoidal and gluing identities by
functoriality. Its boundary trace and Wilson labels agree with the intended
arithmetic ones only under maps establishing those equalities. An arithmetic
decomposition of a closed object could identify its amplitude with an
operator trace. A separate character formula could then identify that trace
with \(\xi(s)\). A modular category alone supplies neither the arithmetic
bordism functor nor the analytic character identity. If the target contains
infinitely many Fock states, the finite trace restriction following
\eqref{v4-arith:cs:eq:c2-modular} must also be addressed.

\subsection{Factorization from fields}

There is a useful positive construction before imposing factorization
descent. Suppose \(U\mapsto\mathcal F(U)\) is a contravariant field-groupoid
functor on opens, finite on the regions in use, and takes finite disjoint
unions to products. Define \(A(U)=\mathbb C^{\pi_0\mathcal F(U)}\). For
disjoint \(U_i\subset V\), set
\begin{equation}
 \mu_{U_i,V}(f_1\otimes\cdots\otimes f_r)(b)
             =\prod_i f_i(b|_{U_i}).
 \label{v4-arith:cs:eq:field-prefactorization}
\end{equation}

Restriction of fields and multiplication of functions prove associativity of
these maps, their symmetry and the empty-input unit. Thus they define a
prefactorization algebra of finite field functions. Weighted integration is
an additional functional on it. Compatibility of that functional with gluing
uses the action and the homotopy-fiber measure.

Descent of the field groupoids does not ensure factorization descent of
these function spaces. A circle gives a counterexample. Take finite
\(C_2\)-bundles and the cover by all complements of one point. It is a Weiss
cover: every finite set of points lies in one member. Each nonempty finite
intersection is a finite union of intervals and has one isomorphism class of
bundles, although it retains automorphisms.

Every function space in its \v{C}ech diagram is therefore \(\mathbb C\),
with identity transition maps. The cover nerve has every finite simplex, so
its homotopy colimit is \(\mathbb C\).

On the circle itself there are two holonomies, and the function space is
\(\mathbb C^2\). The descent map has image the constants and misses the sign
character.

Consequently the arithmetic factorization problem can begin with the maps
\eqref{v4-arith:cs:eq:field-prefactorization}, but must construct or replace
their descent in the intended recipient. Derived or categorical coefficients
can retain data lost by functions on isomorphism classes. Their
factorization maps, support and continuity still need a proof. This issue
already occurs before an archimedean Hilbert extension.

\subsection{Fully dualizable objects}

The framed cobordism hypothesis starts from a symmetric monoidal
\((\infty,3)\)-category and a fully dualizable object. It uses the object's
dual, evaluation, coevaluation and the required adjoints through the
morphism levels, with their coherences. For the Morita three-category of
finite tensor categories, Douglas--Schommer-Pries--Snyder prove that
separable tensor categories are fully dualizable. In characteristic zero
this includes finite semisimple tensor categories. Their theorem gives a
framed three-dimensional theory. An oriented extension uses the appropriate
homotopy fixed-point structure.

A categorical arithmetic recipient with trace and restricted-tensor
conditions must itself be constructed: composition must exist and preserve
those conditions. A monoidal comparison with the ordinary Morita category
must transport the actual duality and adjunction data. A list of candidate
objects and bimodules does not prove these facts. Once they hold, the framed
theory exists in that target. An arithmetic bordism comparison, the boundary
trace map and the Frobenius labels remain separate inputs. Agreement with a
proposed RT theory requires comparison of their extended data in a common
category.

\subsection{A six-functor trace}

Suppose an arithmetic space \(X^\diamond\) has coefficient categories
\(\mathcal C(-)\). Let its structure morphism \(f:X^\diamond\to\mathrm{pt}\)
admit the required functors. For \(B\in\mathcal C(\mathrm{pt})\), the
expression \(f_!f^!B\) then has a definite type. If it is dualizable, its
identity has trace
\[
 \operatorname{Tr}(1_{f_!f^!B})
          \in\operatorname{End}_{\mathcal C(\mathrm{pt})}(1).
\]
A scalar value uses a specified map from this endomorphism ring to the
scalar field. A regularized trace, if needed in place of dualizability,
requires its own domain and normalization.

This produces a trace expression on the supplied geometry. The global
arithmetic diamond and its archimedean boundary are additional inputs. So
are the realization of the Hecke trace as \(B\) and the Euler-product trace
identity. Local geometric Langlands and local diamond geometry provide
arithmetic objects that a comparison may use. They do not by themselves
define all of this global data.

\subsection{Comparing Hilbert realizations}

A common positive pairing does give a strong comparison once the state
spaces have been constructed. Let \(E_\infty\) be a specified complex trace
realization, with a positive semidefinite Hermitian form \(q_\infty\). Write
its nullspace as \(N_q\).

\begin{proposition}[Boundary Hilbert transport]
\label{v4-arith:cs:prop:hilbert-transport}
Suppose \(I_B:E_\infty\to H_B\) and \(I_P:E_\infty\to H_P\) have dense
ranges and satisfy
\[
 \langle I_Bv,I_Bw\rangle
   =q_\infty(v,w)
   =\langle I_Pv,I_Pw\rangle.
\]
There is a unique unitary \(T:H_B\to H_P\) with \(TI_B=I_P\). If closed
operators \(A_B,A_P\) have cores \(I_B\mathcal D\), \(I_P\mathcal D\),
respectively, and \(TA_BI_Bv=A_PI_Pv\) for \(v\in\mathcal D\), then
\[
 T\operatorname{Dom}A_B
        =\operatorname{Dom}A_P,
 \qquad TA_B=A_PT.
\]
For normal operators this also intertwines their bounded spectral functions.
\end{proposition}

\begin{proof}
Cauchy--Schwarz makes \(N_q\) orthogonal to all of \(E_\infty\). Both
realization kernels equal \(N_q\), so \(I_Bv\mapsto I_Pv\) is a well-defined
isometry on dense subspaces. It extends uniquely onto \(H_P\). The graph of
each closed operator is the closure of its graph on the specified core.
Applying \(T\oplus T\) gives the domain identity and the full intertwining
relation. The normal functional calculus then proves the last assertion.
\end{proof}

Boundary operations must preserve \(N_q\) to descend even algebraically.
Their closed generators need the domains in this proposition. A further
identification with logarithmic scaling on \(L^2(\mathbb R_{>0},dx/x)\)
requires another unitary with its domain identity. Polynomial Fock energy
has eigenvectors, whereas the full-line logarithmic translation generator
has none. Those representations cannot be identified by a unitary preserving
these generators.

Positive Hilbert descent is also distinct from the complete Weil pairing.
Let \(\mathcal T\) be the chosen Weil test algebra with involution \(*\).
Its complete signed functional \(W\) includes the polar, archimedean and
prime terms.

If a map \(J:\mathcal T\to E_\infty\) satisfies
\begin{equation}
 q_\infty(Jf,Jg)=W(f*g^*)
                \qquad(f,g\in\mathcal T),
 \label{v4-arith:cs:eq:weil-comparison}
\end{equation}
then \(W(f*f^*)\ge0\). This follows by setting \(g=f\). The substantive
assertion is \eqref{v4-arith:cs:eq:weil-comparison} on the full required
test domain, with its completion and normalization. A positive form on an
auxiliary finite gauge or oscillator space does not supply it. Conversely,
positivity of the full Weil form does not construct the arithmetic bordism
functor.

\section{Other physical realizations and the remaining maps}
\label{v4-arith:cs:remaining}

Arithmetic Chern--Simons is one possible physical formulation. It does not
exclude others. The class-\(\mathcal S\) construction compactifies the
six-dimensional \((2,0)\) theory on a complex curve, with its twist and
defect data, to obtain four-dimensional \(\mathcal N=2\) theories.

For an unpunctured complex curve its Hitchin base uses
\(\bigoplus_iH^0(C,K_C^{d_i})\). Punctures change the allowed poles.
Replacing \(K_C\) by a metrized arithmetic dualizing line does not specify a
Higgs field, Hitchin equations or a six-dimensional compactification. Krull
dimension cannot be added to a real manifold dimension to prove that every
arithmetic version is impossible. The question is to construct the geometry,
fields and comparison that would make such a version defined.

Likewise, a local Hecke character is a label, not the one-dimensional defect
carrying it. Local class field theory relates its appropriate finite or Weil
realization to local arithmetic data. Turning it into a boundary module
requires a functor to the specified defect category, including its ramified
action. The local character space can itself have positive-dimensional
families. Choosing a point in it does not assign a dimension to the physical
defect.

The protected chiral construction of
Beem--Lemos--Liendo--Peelaers--Rastelli--van Rees starts with an actual
four-dimensional superconformal theory and its protected cohomology. Its
algebra and Schur trace retain those fields and symmetries. An arithmetic
Euler trace has a simpler independent construction: on the finite occupation
states with energies \(\sum_pn_p\log p\), the trace of \(e^{-sH}\) is
\(\prod_p(1-p^{-s})^{-1}\) for \(\operatorname{Re}s>1\). This does not
identify the space with protected cohomology. Frobenius is an endomorphism,
not automatically an integer grading. A symbol \(t^{\operatorname{Frob}}\)
requires a defined functional calculus or a specified graded realization.

The arithmetic comparison therefore remains an object and map problem. The
required data are the arithmetic bulk geometry, field action, coherent
measures and gluing. One must construct the boundary trace realization, its
pairing and the Frobenius Wilson comparison. The partition trace also
requires its analytic domain.

The finite torsors, abelian linking correlations, Haar realization and free
physical period map give positive constructions in four specified settings.
Their extension to one arithmetic theory requires the maps just identified.
The full Weil comparison adds a further identity on its complete test
algebra.

\section*{Sources}

\begin{itemize}
\item D.~Freed and F.~Quinn,
\emph{Chern--Simons theory with finite gauge group}. \emph{Communications in
Mathematical Physics} 156 (1993), 435--472. Theorems~1.7 and~2.13 and
Lemma~2.4 give the finite topological theory.
\url{https://arxiv.org/abs/hep-th/9111004}.
\item M.~Kim, \emph{Arithmetic Chern--Simons theory I}, \S\S1--2.
The finite arithmetic action is defined there.
\url{https://arxiv.org/abs/1510.05818}.
\item H.-J.~Chung, D.~Kim, M.~Kim, J.~Park and H.~Yoo.
\emph{Arithmetic Chern--Simons theory II}, \S\S2--4 and Appendix~B. These
contain the boundary torsor, compact-coefficient construction, decomposition
formula and gauge prisms. \url{https://arxiv.org/abs/1609.03012}.
\item H.-J.~Chung, D.~Kim, M.~Kim, G.~Pappas, J.~Park and H.~Yoo.
\emph{Abelian arithmetic Chern--Simons theory and arithmetic linking
numbers}. Theorems~1.1 and~4.1 and \S3 give the linking pairing, Gaussian
correlation, relative construction and class invariant.
\url{https://arxiv.org/abs/1706.03336}.
\item C.~Douglas, C.~Schommer-Pries and N.~Snyder,
\emph{Dualizable tensor categories}. See Theorem~3.4.3 and
Corollaries~3.4.6--3.4.8 for the framed tensor-category field theory.
\url{https://arxiv.org/abs/1312.7188}.
\item The Stacks Project, Theorem~10.34.1,
\url{https://stacks.math.columbia.edu/tag/00FV}. This is the Nullstellensatz
used in the finite-orbit argument.
\item D.~Gaiotto, \emph{\(\mathcal N=2\) dualities},
\url{https://arxiv.org/abs/0904.2715}. For the class-\(\mathcal S\) and
protected chiral constructions, see also
Beem--Lemos--Liendo--Peelaers--Rastelli--van Rees, \emph{Infinite chiral
symmetry in four dimensions}, \S6.1, \url{https://arxiv.org/abs/1312.5344}.
\end{itemize}
